\documentclass[a5paper, 10pt]{article}


\usepackage{cite, amsmath, amssymb, booktabs, lastpage, graphicx, float}
\usepackage[hidelinks]{hyperref}

\usepackage{geometry}
\geometry{margin=1.5cm,a5paper,twoside,inner=2cm}

\usepackage{setspace}


\usepackage{amsthm}

\theoremstyle{plain}
\newtheorem{lemma}{Lemma}
\newtheorem{theorem}{Theorem}
\newtheorem{conjecture}{Conjecture}
\newtheorem{proposition}[theorem]{Proposition}
\newtheorem*{corollary}{Corollary}
\renewcommand{\qedsymbol}{\rule{.7em}{.7em}}

\theoremstyle{remark}
\newtheorem*{remark}{Remark}
\newtheorem*{note}{Note}
\newtheorem{case}{Case}

\theoremstyle{definition}
\newtheorem{definition}{Definition}
\newtheorem{example}{Example}

%% Changing the title, author, etc.
\usepackage{graphicx}
\makeatletter
\renewcommand{\maketitle}{
    \begin{center}
    \rule[.2em]{\textwidth}{0.353mm}
        \begin{minipage}[m]{0.35\textwidth}
            {\scriptsize
                \begin{center}
                    \textsf{\textbf{\huge MATCH}\\
                        \textit{Communications in Mathematical\\
                            and in Computer Chemistry}
                    }
            \end{center}}
        \end{minipage}\hfill
        \begin{minipage}[m]{0.65\textwidth}
            \begin{flushright}
                \baselineskip=10pt
                {\scriptsize\sffamily{\itshape  MATCH Commun. Math. Comput. Chem.}
                \textbf{\vol{}}
                (\pubyear)
                    \thepage--\pageref{LastPage}}\\
                {\scriptsize\sffamily {\bfseries ISSN:} 0340--6253}\\
                {\scriptsize\sffamily \textbf{doi:} \doi{10.46793/match.\vol-\num.\id}}
            \end{flushright}
        \end{minipage}
        \rule[1em]{\textwidth}{.353mm}
        \baselineskip=0.30in
        {\Large\bfseries \@title} \par
        \vspace{5mm}
        \baselineskip=0.2in
        {\large\bfseries \@author}\par
        \vspace{1mm}
        {\it \@address} \par
        {\small\tt \@email} \par
        \vspace{3mm}
        {\small (Received \@date)} \par
    \end{center}
    \vspace{3mm}
}
\newcommand{\address}[1]{\def\@address{#1}}
\newcommand{\email}[1]{\def\@email{#1}}
\newcommand{\doi}[1]{\href{https://doi.org/#1}{#1}}

\newcommand{\vol}{00}
\newcommand{\num}{00}
\newcommand{\pubyear}{0000}
\newcommand{\id}{xxxxx}
%%setting page-number
\setcounter{page}{1}


\thispagestyle{empty}
\makeatother

%% Making << \acknowledgment >> command
\newcommand{\acknowledgment}[1]{\vspace{5mm}\singlespacing
    {\noindent\textbf{\textit{Acknowledgment\/}:} #1}
}

%% Setting page numbers
\usepackage{fancyhdr}

%% now reset the headers and footers
\fancyhead{}
\fancyfoot{}
\renewcommand{\headrulewidth}{.3pt}
\voffset = 18pt
\headsep = 3pt

\renewcommand*{\thefootnote}{\fnsymbol{footnote}}

\usepackage[margin=1cm,%
            font=footnotesize,%
            format=hang,%
            labelsep=period,%
            labelfont=bf]{caption}

%%%%%%%%%%%%%%%%%%%%%%%%%%%%%%%%%%%%%%%%%%%%%%%%%%
% PLACE ADDITIONAL PACKAGES HERE (if it is needed):
\usepackage[T1]{fontenc}
\usepackage{lmodern}
\usepackage{mathtools}
\usepackage{array}
\usepackage{enumitem}
\setlist[itemize]{topsep=2pt,itemsep=2pt}
\setlist[enumerate]{topsep=2pt,itemsep=2pt}
\emergencystretch=3em
\hbadness=2000
\setlength{\marginparwidth}{2cm}
%%%%%%%%%%%%%%%%%%%%%%%%%%%%%%%%%%%%%%%%%%%%%%%%%%

%% make the odd pages have the section name on the top right
\fancyhead[RO]{\footnotesize \thepage}
%% make the even pages have the chapter name on the top left
\fancyhead[LE]{\footnotesize \thepage}

%% making page-numbers visible
\pagestyle{fancy}

% Custom macros for the Loyola index family
    \newcommand{\Wuzi}{LO(G;\alpha,\beta,\gamma)}
  \newcommand{\edgepsi}{\psi(x,y;\alpha,\beta,\gamma)}
  \newcommand{\bid}{\mathcal{F}_{\mathrm{BID}}}
% Path to figures
\graphicspath{{figures/}{./}{../figures/}}

% robust figure inclusion: render if present (root or figures/), else a small note
\newcommand{\safefig}[2]{%
  \IfFileExists{#1}{\includegraphics[width=#2\linewidth]{#1}}{%
   \IfFileExists{figures/#1}{\includegraphics[width=#2\linewidth]{#1}}{%
    \fbox{\parbox[c][3cm][c]{0.8\linewidth}{\centering Figure not in project --- upload \texttt{\detokenize{#1}}.}}}}%
}

%%%%%%%%%%%%%%%%%%%%%%%%%%%%%%%%%%%%%%%%%%%%%%%%%%%%%%%%%%%%%%%%%%%%%%%%%%%%%%%%%%%%%%%

\title{The Loyola Index Family, Extremal Graph Characterizations and Sensitivity Analysis}
\author{Advik Natarajan$^{a}$, Ganesh Shiwakoti$^{b}$, Michael Arockiaraj$^{a}$\footnote{Corresponding author.}}

\address{$^a$Department of Mathematics, Loyola College, Chennai, India\\
    $^b$Department of Mathematics and Statistics, Florida Atlantic University, USA
    }

\email{marockiaraj@gmail.com}

\date{\today}

% !TeX program = pdflatex

\begin{document}

\maketitle

\begin{abstract}
  The two-parameter Gutman--Milovanovi\'c family
  $M_{\alpha,\beta}(G) = \sum_{uv \in E(G)} (d_u d_v)^{\alpha} (d_u + d_v)^{\beta}$
  unifies the Zagreb, Randi\'c, sum-connectivity and related degree-based
  indices. We study its normalized degree-imbalance exponential
  extension, the \emph{Loyola index family}
  \begin{equation*}
  LO(G;\alpha,\beta,\gamma)
  \;=\; \sum_{uv \in E(G)}
  (d_u d_v)^{\alpha}\,(d_u + d_v)^{\beta}\,
  \exp\!\left(\gamma\,\frac{|d_u - d_v|}{d_u + d_v}\right),
  \end{equation*}
  with $LO(G;\alpha,\beta,0) = M_{\alpha,\beta}(G)$ identically.

  We prove that the imbalance direction is affinely independent of the
  two parent directions on the admissible degree-pair set of graphs with
  maximum degree at least $3$ (a class-level, not per-graph, statement), that
  $\log LO$ is globally convex in $(\alpha,\beta,\gamma)$ with gradient
  and Hessian given by the mean and covariance of the edge-feature
  vector under an induced edge measure, and we extend two-sided bounds
  from the parent family to $\gamma \neq 0$. On molecular graphs of
  maximum degree at most four the entire family is confined to the
  finite-dimensional space of edge-degree-pair counts; we quantify this
  ceiling exactly, exhibiting octane-isomer pairs that no
  vertex-degree-based index can distinguish, and showing that the
  pure-$\gamma$ points attain the maximal degree-pair discrimination on
  trees of orders $7$--$12$. A matched-budget nested leave-one-out
  comparison against the optimized parent family on five octane
  properties, repeated over one hundred random-search seeds at two
  budgets, finds that the free imbalance parameter usually improves
  out-of-sample prediction on one property, usually degrades it on
  three, and is mixed on the fifth, with neither tendency sign-invariant across candidate
  streams, so $\gamma$ is presented as a structurally interpretable
  and provably independent direction rather than a predictive
  upgrade;
  under the published structure-sensitivity protocol of Furtula,
  Gutman and Dehmer, the pure-$\gamma$ points show the largest
  sensitivity-to-abruptness ratios in the tested comparison set.
  \end{abstract}

\onehalfspacing

% =====================================================================
\section{Introduction}\label{sec:intro}
  % =====================================================================

  The literature on degree-based topological indices in chemical
  graph theory has grown substantially since the seminal Wiener
  index, the Randi\'c index,
  and the Zagreb indices of Gutman and
  Trinajsti\'c.
  Among the many parametric extensions are the general Randi\'c
  index $R_{\alpha}$~\cite{BollobasErdos1998}, the general
  sum-connectivity index
  $\chi_{\alpha}$~\cite{ZhouTrinajstic2010}, the forgotten
  index $F$ revisited by Furtula and
  Gutman~\cite{FurtulaGutman2015F}, the geometric--arithmetic index
  $GA$~\cite{VukicevicFurtula2009GA}, the Albertson irregularity
  index~\cite{Albertson1997} and the inverse $\sigma$-index
  problem~\cite{GutmanToganYurttasCevikCangul2018Sigma}.

  While historical literature focused primarily on algebraic formulations
  to capture molecular branching and structural invariants, a distinct
  shift toward geometric interpretations occurred when Gutman introduced
  the Sombor index family \cite{Gutman2021Sombor, Gutman2021SomborBasic}.
  This formulation led to work such as establishing fundamental bounds
  \cite{MilovanovicMilovanovicMatejic2021Sombor} and evaluating how these
  new geometric invariants correlated with traditional degree-based metrics
  \cite{WangMaoLiFurtula2022SomborRelations}. Much of this initial work
  concentrated on identifying extremal bounds within fundamental graph
  classes, particularly optimizing the general and standard variants of the
  index over molecular trees and related networks
  \cite{DasCevikCangulShang2021Sombor, DasGutman2022SomborTrees,
  AhmadFarooqDas2025Extremal}.

   Contemporary studies have introduced modified metrics such as the mean
  Sombor index \cite{MendezBermudezAguilarSanchez2022MeanSombor}, the
  elliptic Sombor index \cite{GutmanFurtulaOz2024Elliptic}, and the
  hyperbolic Sombor index \cite{BarmanDas2026HSO}. Concurrently, parameters
  like the diminished Sombor index \cite{Movahedi2025DSO,
  MovahediGutmanRedzepovicFurtula2026} and the hybrid vertex-edge KG-Sombor
  index \cite{KosariDehgardiKhan2023Lower} have been developed to handle
  unique structural constraints. Beyond purely theoretical analysis, these
  descriptors are increasingly integrated into practical quantitative
  structure-property relationship (QSPR) models and spectral evaluations.
  For instance, hybrid geometric-harmonic-Zagreb descriptors are now
  utilized to optimize the predictive analysis of polycyclic aromatic
  hydrocarbons \cite{ArockiarajPaulClement2023Hybrid}, while modified
  reverse degree invariants help characterize the graph energy and entropy
  of complex 2D metal-organic frameworks \cite{KalaamGreeniArockiaraj2024}.
  This rapid diversification of structural metrics is currently being
  unified through generalized matrix methods
  \cite{ZuoRatherImranAli2022Modified}, dual coindex formulations
  \cite{GujiWali2025Lanzhou}, simplicial network assessments
  \cite{Shang2022Simplicial}, and finite-dimensional linear frameworks
  \cite{Rada2026LinearVDB}, as documented in recent comprehensive
  literature surveys \cite{Gutman2026SomborSurvey}.

  The natural two-parameter unification of the product-based and
  sum-based degree indices is the \emph{Gutman--Milovanovi\'c family},
  introduced by Gutman, Milovanovi\'c and
  Milovanovi\'c~\cite{GutmanMilovanovic2020AKCE} and studied
  intensively in the recent bounds and QSPR literature
  \cite{GranadosJMC2025GM, MolinaSigarreta2025GM, PortillaJMC2025GM}:
  \begin{equation*}
  M_{\alpha,\beta}(G) \;=\; \sum_{uv \in E(G)}
  (d_u d_v)^{\alpha} (d_u + d_v)^{\beta}, \qquad
  \alpha, \beta \in \mathbb{R}.
  \end{equation*}
  Its special cases are catalogued
  in~\cite{MolinaSigarreta2025GM}: $M_{0,1} = M_1$ (first Zagreb),
  $M_{1,0} = M_2$ (second Zagreb), $M_{0,2} = HM$ (hyper-Zagreb),
  $M_{-1,0} = {}^{m}\!M_2$ (modified second Zagreb),
  $M_{-1/2,0} = R$ (Randi\'c),
  $M_{0,-1/2} = \chi$ (sum-connectivity~\cite{ZhouTrinajstic2009sumconn}),
  $2M_{0,-1} = H$ (harmonic), $M_{1,-1} = ISI$ (inverse sum indeg),
  $2M_{1/2,-1} = GA$ (geometric-arithmetic),
  $\tfrac12 M_{-1/2,1} = AG$ (arithmetic-geometric), and the whole
  general Randi\'c family $R_\alpha = M_{\alpha,0}$
  \cite{BollobasErdos1998} and general sum-connectivity family
  $\chi_\beta = M_{0,\beta}$ \cite{ZhouTrinajstic2010} as parameter
  axes. We claim no novelty for this backbone, for its reduction
  catalogue, or for its generic tunability: those belong to the
  Gutman--Milovanovi\'c literature.

  In this paper we study a three-parameter extension of
  $M_{\alpha,\beta}$, the \emph{Loyola index family}, obtained by
  attaching a separately parameterized exponential factor in the
  normalized degree imbalance
  $q_{uv} = |d_u - d_v|/(d_u + d_v)$ of each edge:
  $LO(G;\alpha,\beta,\gamma) = \sum_{uv}
  (d_u d_v)^{\alpha}(d_u+d_v)^{\beta} e^{\gamma q_{uv}}$, so that
  $LO(G;\alpha,\beta,0) = M_{\alpha,\beta}(G)$ identically.
  Positive $\gamma$ exponentially upweights degree-imbalanced edges,
  negative $\gamma$ downweights them, $q_{uv} = 0$ on equal-degree
  edges, and on regular graphs $\gamma$ has no effect. The questions
  this paper answers are therefore about $\gamma$ alone: is the
  imbalance direction mathematically independent of the parent
  parameters, what analytic structure does it carry, and does it add
  measurable structural or predictive information beyond an optimized
  Gutman--Milovanovi\'c parent?

  On the mathematical side, we prove that the answer to independence
  is affirmative on the admissible degree-pair set of graphs of maximum
  degree at least three (Proposition~\ref{prop:indep}; on an individual
  graph $\gamma$ can be absorbed by the parent parameters,
  Remark~\ref{rem:ident}), establish global log-convexity of
  the family in $(\alpha,\beta,\gamma)$ with an exact
  mean/covariance interpretation of its derivatives
  (Theorem~\ref{thm:logconvex}), give closed-form values on standard
  graph classes, and derive two-sided bounds that extend inequalities
  known for $M_{\alpha,\beta}$
  \cite{MolinaSigarreta2025GM, GranadosJMC2025GM} to the
  $\gamma \neq 0$ regime, several sharp in their equality cases. On
  the chemometric side we quantify the exact information ceiling of
  the entire vertex-degree-based class on molecular graphs, run a
  matched-budget nested comparison of $LO$ against its optimized
  parent, and evaluate the family under the published
  structure-sensitivity protocol of Furtula, Gutman and
  Dehmer~\cite{FurtulaGutmanDehmer2013SS}; no universal ranking or
  predictive-superiority claim is made.

  The distinguishing feature of the Loyola family within the
  degree-based BID class is the multiplicative
  $\exp(\gamma\, q_{uv})$ factor, where
  $q_{uv} = |d_u - d_v|/(d_u + d_v) \in [0, 1)$ is the
  normalized degree imbalance. Writing the summand as
  $\exp\bigl(\alpha \ln(d_u d_v) + \beta \ln(d_u + d_v) +
  \gamma\, q_{uv}\bigr)$ shows that the family is a
  three-parameter subfamily of the exponential vertex-degree-based
  indices introduced by Rada~\cite{Rada2019ExpVDB} and studied further
  in extremal settings by Sigarreta~\cite{Sigarreta2022ExpVDB}; general
  sufficient conditions for tree-extremal behaviour of such exponential
  edge-weight sums are known~\cite{HuMATCH2022, GaoAMC2024}, and no
  extremal-ordering novelty is claimed here. The published
  one-parameter exponential indices place their parameter on a single
  classical generator inside the exponent; the scoped contribution of
  this family is the multi-parameter \emph{linear combination} in the
  exponent in which the two Gutman--Milovanovi\'c basis terms act as
  the factors $(d_u d_v)^{\alpha}(d_u + d_v)^{\beta}$ while the
  imbalance coordinate carries its own parameter $\gamma$. The same
  imbalance ratio appears in the diminished Sombor
  index~\cite{MovahediGutmanRedzepovicFurtula2026} and in
  irregularity-based BID variants; to our knowledge, and within the
  scope of the literature search documented in the reproducibility
  package, no published index couples it multiplicatively, with its
  own parameter, to the Gutman--Milovanovi\'c family. Since
  $0 \le q_{uv} < 1$, the factor lies in $[1, e^{\gamma})$
  for $\gamma > 0$, equals $1$ for $\gamma = 0$, and lies in
  $(e^{\gamma}, 1]$ for $\gamma < 0$; it reduces to $1$ on regular
  graphs. The exponential
  form is a structural choice rather than one derived from a physical
  model: it lets $\gamma$ act multiplicatively and only on irregularity,
  vanishing on regular graphs.

  Also, it is useful to demonstrate the descriptor-space ceiling in
  which the entire family lives. On hydrogen-suppressed molecular
  graphs of maximum degree at most $4$, every BID index is a linear
  functional of the ten unordered edge-degree-pair counts
  $m_{ij}(G) = |\{uv \in E(G) : \{d_u, d_v\} = \{i, j\}\}|$ for
  $1 \le i \le j \le 4$; this finite-dimensional view of
  vertex-degree-based indices has recently been placed on a formal
  linear-algebraic footing by Rada~\cite{Rada2026LinearVDB}. Thus
  the Loyola family, like every BID family on this molecular class,
  lies inside a fixed ten-dimensional descriptor
  space, no matter how
  its three real parameters are chosen.

  In sum, the contributions of this paper are: (i) the affine
  independence of the imbalance direction from the two
  Gutman--Milovanovi\'c directions on the admissible degree-pair set
  for $\Delta \ge 3$ (Proposition~\ref{prop:indep}); (ii) global log-convexity of the
  family with exact mean/covariance formulas for its parameter
  derivatives (Theorem~\ref{thm:logconvex}); (iii) two-sided bounds
  extending parent-family inequalities to $\gamma \neq 0$, several
  sharp; (iv) an exact quantification of the information ceiling of
  the whole vertex-degree-based class on molecular graphs, including
  octane-isomer pairs indistinguishable by every such index
  (Proposition~\ref{prop:collision}) and floor-attainment of the
  pure-$\gamma$ points across tree orders $7$--$12$; and (v) an
  honest matched-budget nested assessment of whether $\gamma$ adds
  out-of-sample predictive information beyond the optimized parent
  (in a $100$-seed analysis it usually helps on one of five octane
  properties, usually hurts on three, and is mixed on the fifth, with
  neither tendency
  sign-invariant across candidate streams), together
  with the family's behaviour under the published
  structure-sensitivity protocol~\cite{FurtulaGutmanDehmer2013SS},
  validated against published decane-class
  values~\cite{BarmanDas2026HSO}.
  The family is thus offered not as a universally superior predictor
  but as a controlled, analytically transparent probe of the BID
  descriptor space whose third direction is provably independent of
  its established parent.

% =====================================================================
\section{Preliminaries}\label{sec:prelim}
% =====================================================================

\subsection{Notation}\label{sec:notation}

Let $G = (V, E)$ be a finite simple connected graph of order
$n = |V|$ and size $m = |E|$. For a vertex $v \in V$, the degree
$d_v = d_G(v)$ counts the edges incident to $v$. The maximum and
minimum degrees are denoted $\Delta = \Delta(G)$ and
$\delta = \delta(G)$. In the chemometric sections we work with
hydrogen-suppressed molecular graphs, where the assumption
$\Delta \le 4$ is natural for organic molecules built from
$\{$C, N, O, S, P, halogens$\}$.

The classical graph families used throughout are: the complete
graph $K_n$, the cycle $C_n$, the path $P_n$, the star $S_n$, the complete
bipartite graph $K_{p,q}$, the hypercube $Q_k$ (the
$k$-dimensional binary hypercube; $2^k$ vertices,
$k \cdot 2^{k - 1}$ edges, $k$-regular), the wheel $W_n$ (an
$(n - 1)$-cycle plus a universal vertex), the friendship graph
$F_p$ ($p$ triangles sharing a common vertex), and the
$r$-regular graphs.

For an edge $e = uv \in E(G)$ we write $(d_u, d_v)$ for its
\emph{degree pair} understood as an unordered multiset
$\{d_u, d_v\}$.

\subsection{Definition}\label{sec:def}

\begin{definition}[Loyola index family]\label{def:wuzi}
Let $G = (V, E)$ be a simple connected graph. For
$\alpha, \beta, \gamma \in \mathbb{R}$, the \emph{Loyola index family}
is defined as,
\begin{equation}\label{eq:wuzi}
\Wuzi
\;=\; \sum_{uv \in E(G)}
(d_u d_v)^{\alpha}\,(d_u + d_v)^{\beta}\,
\exp\!\left(\gamma\,\frac{|d_u - d_v|}{d_u + d_v}\right)
\end{equation}
\end{definition}

Equation~\eqref{eq:wuzi} is well defined for all
$\alpha, \beta \in \mathbb{R}$ on graphs without isolated vertices
(so that $d_u, d_v \ge 1$ on every edge). The expression is
invariant under permutations of $(u, v)$ because
$d_u d_v = d_v d_u$, $d_u + d_v = d_v + d_u$, and
$|d_u - d_v| = |d_v - d_u|$; hence $\Wuzi$ depends only on the
multiset of unordered edge degree pairs of~$G$. When $\gamma = 0$
the family reduces identically to the Gutman--Milovanovi\'c
family~\cite{GutmanMilovanovic2020AKCE, MolinaSigarreta2025GM},
\begin{equation*}
LO(G;\alpha,\beta,0) \;=\; M_{\alpha,\beta}(G)
\;=\; \sum_{uv \in E(G)} (d_u d_v)^{\alpha} (d_u + d_v)^{\beta},
\end{equation*}
whose classical special cases are catalogued in the Introduction.
Throughout, $q_{uv} = \mathrm{imb}(d_u, d_v)$, where
$\mathrm{imb}(x, y) = |x - y|/(x + y)$, denotes the
normalized degree imbalance of the edge $uv$; positive $\gamma$
exponentially upweights degree-imbalanced edges, negative $\gamma$
downweights them, and on regular graphs every $q_{uv} = 0$ so that
$\gamma$ has no effect and
$LO(G;\alpha,\beta,\gamma) = M_{\alpha,\beta}(G)$ for all $\gamma$.
\subsection{Closed-form values on standard graphs}\label{sec:closedforms}

  \begin{proposition}\label{prop:closedforms}
  Let $\alpha, \beta, \gamma \in \mathbb{R}$. Then:
  \begin{enumerate}
  \item  $LO(K_n; \alpha, \beta, \gamma) =
    \dfrac{n(n-1)}{2}\,2^{\beta}(n - 1)^{2\alpha + \beta}$ for $n \ge 2$.
  \item  $LO(C_n; \alpha, \beta, \gamma) = n \cdot 4^{\alpha +
  \beta}$
    for $n \ge 3$.
  \item
    $LO(P_n; \alpha, \beta, \gamma) =
    \begin{cases}
      2^{\beta}, & n = 2,\\[2pt]
      2 \cdot 2^{\alpha} 3^{\beta} e^{\gamma/3}, & n = 3,\\[2pt]
      2 \cdot 2^{\alpha} 3^{\beta} e^{\gamma/3} + (n - 3) \cdot 4^{\alpha +
  \beta},
         & n \ge 4.
    \end{cases}$
  \item  $LO(S_n; \alpha, \beta, \gamma) =
    (n - 1)\,(n - 1)^{\alpha} n^{\beta} \exp\!\left(\gamma\,\dfrac{n -
  2}{n}\right)$
    for $n \ge 2$.
  \item  $LO(K_{p,q}; \alpha, \beta, \gamma) =
    pq\,(pq)^{\alpha}(p + q)^{\beta}\exp\!\left(\gamma\,\dfrac{|p - q|}{p +
  q}\right)$.
  \item  $LO(Q_k; \alpha, \beta, \gamma) =
    k \cdot 2^{k - 1} \cdot 2^{\beta} k^{2\alpha + \beta}$ for $k \ge 1$.
  \item  For $n \ge 5$, $W_n$ has $n - 1$ spoke edges with
    degree pair $(3, n - 1)$ and $n - 1$ rim edges with degree pair
    $(3, 3)$, so
    \[
    LO(W_n; \alpha, \beta, \gamma)
    = (n - 1)\Bigl[\,3^{\alpha} (n - 1)^{\alpha}(n + 2)^{\beta}
         \exp\!\left(\gamma\,\tfrac{n - 4}{n + 2}\right)
       + 9^{\alpha} 6^{\beta}\Bigr].
    \]
  \item  For $p \ge 1$, $F_p$ has $2p$ spoke edges
    $(2, 2p)$ and $p$ outer edges $(2, 2)$, so
    \[
    LO(F_p; \alpha, \beta, \gamma)
    = p \cdot 4^{\alpha + \beta}
     + 2p\,(4p)^{\alpha}(2p + 2)^{\beta}
         \exp\!\left(\gamma\,\tfrac{p - 1}{p + 1}\right).
    \]
  \item  For an $r$-regular graph $G$ with $m$ edges,
    $\Wuzi = m \cdot 2^{\beta} r^{2\alpha + \beta}$.
  \end{enumerate}
  \end{proposition}

  \begin{proof}
  Each case is obtained by enumerating the edge degree pairs of the
  graph and applying Definition~\ref{def:wuzi}. For $r$-regular
  graphs and for $K_n, C_n, Q_k$ (all regular), every edge has
  $d_u = d_v = r$, so $\mathrm{imb} = 0$ and
  $(d_u d_v)^{\alpha}(d_u + d_v)^{\beta}
     = r^{2\alpha}(2r)^{\beta} = 2^{\beta} r^{2\alpha + \beta}$.
  Items (1), (2), (6), and (9) follow immediately. For $S_n$ all
  $n - 1$ edges have degree pair $(1, n - 1)$ with
  $\mathrm{imb} = (n - 2)/n$; item (4) follows. For $K_{p, q}$ all
  $pq$ edges have degree pair $(p, q)$; item (5) follows. For $P_n$
  ($n \ge 4$) there are exactly two pendant edges with pair $(1, 2)$
  and $n - 3$ interior edges with pair $(2, 2)$; item (3) follows by
  direct summation. The boundary cases $n = 2, 3$ are immediate.
  Items (7) and (8) are obtained by partitioning the edges of $W_n$
  and $F_p$ into their two edge-degree-pair classes respectively and
  summing.
  \end{proof}

 % =====================================================================
  \section{Parameter geometry and the role of the imbalance
  parameter}\label{sec:geometry}

  This section establishes the two structural facts that justify
  introducing the third parameter at all: the imbalance direction
  cannot be reproduced by the parent parameters, and the family has a
  globally convex logarithm whose derivatives have an exact
  probabilistic meaning. Throughout, write
  $z_{uv} = \bigl(\ln(d_u d_v),\, \ln(d_u + d_v),\, q_{uv}\bigr)
  \in \mathbb{R}^3$ for the feature vector of the edge $uv$ and
  $\theta = (\alpha, \beta, \gamma)$, so that
  $LO_G(\theta) = \sum_{uv \in E(G)} e^{\theta^{\top} z_{uv}}$.

  \begin{proposition}[Affine independence of the imbalance
  direction]\label{prop:indep}
  Let $\Delta \ge 3$ and let
  $P_{\Delta} = \{(i,j) : 1 \le i \le j \le \Delta\}$ be the
  admissible degree pairs. Then the augmented feature vectors
  $\bigl(1, \ln(ij), \ln(i+j), q_{ij}\bigr)$, $(i,j) \in P_{\Delta}$,
  span $\mathbb{R}^4$. Consequently, if two parameter triples
  $\theta, \theta'$ satisfy
  $e^{\theta^{\top} z} = c\, e^{\theta'^{\top} z}$ for some constant
  $c > 0$ and all $(i,j) \in P_{\Delta}$, then $\theta = \theta'$ and
  $c = 1$; in particular no choice of $(\alpha', \beta')$ and scaling
  constant reproduces the action of $\gamma \neq 0$ on all admissible
  degree pairs. For $\Delta = 2$ the span has dimension $3$, and for
  $\Delta = 1$ (the single pair $(1,1)$, augmented vector
  $(1, 0, \ln 2, 0)$) it has dimension $1$; in both cases the
  conclusion fails.
  \end{proposition}

  \begin{proof}
  It suffices to show that $q_{ij}$ is not an affine function of
  $\bigl(\ln(ij), \ln(i+j)\bigr)$ on $P_{\Delta}$; the four pairs
  $(1,1), (2,2), (1,2), (1,3)$, all admissible once $\Delta \ge 3$,
  witness this. Suppose
  $q_{ij} = a + b \ln(ij) + c \ln(i+j)$ on these pairs. From
  $(1,1)$ and $(2,2)$, where $q = 0$: $a + c \ln 2 = 0$ and
  $a + (2b + 2c) \ln 2 = 0$, giving $a = -c \ln 2$ and $b = -c/2$.
  Substituting into the $(1,2)$ equation
  $\tfrac13 = a + b \ln 2 + c \ln 3$ yields
  $c = \tfrac13 \bigl(\ln 3 - \tfrac32 \ln 2\bigr)^{-1}$, and then the
  $(1,3)$ equation would require
  $\tfrac12 = c\,(\ln 2 - \tfrac12 \ln 3) \approx 0.814$, a
  contradiction. Hence the augmented vectors have rank $4$. If
  $e^{\theta^{\top} z} = c\, e^{\theta'^{\top} z}$ on $P_{\Delta}$,
  then $(\theta - \theta')^{\top} z_{ij} = \ln c$ for all admissible
  pairs, i.e.\ $(\ln c, \theta' - \theta)$ is orthogonal to every
  augmented vector, so $\theta = \theta'$ and $c = 1$. For
  $\Delta = 2$ only the pairs $(1,1), (1,2), (2,2)$ exist and a
  direct computation shows the augmented vectors have rank $3$; for
  $\Delta = 1$ the single augmented vector has rank $1$.
  \end{proof}

  \begin{remark}[Independence versus per-graph
  identifiability]\label{rem:ident}
  Proposition~\ref{prop:indep} is a statement about the edge-weight
  \emph{function} on the full admissible pair set, not about every
  individual graph: on a single graph $G$, the parameters are
  identifiable (up to the scaling noted above) only when the degree
  pairs actually realised by $G$ carry four affinely independent
  augmented feature vectors. On a regular graph nothing about
  $\theta$ is identifiable beyond one linear functional; on a graph
  whose edges realise at most three affinely independent augmented
  vectors, some direction of $\theta$ --- possibly the $\gamma$
  direction --- is unidentifiable from $LO(G;\cdot)$ alone. The
  proposition guarantees that no such degeneracy is forced on the
  whole molecular class $\Delta \ge 3$: the family separates
  parameters as a family of descriptors, even though particular
  graphs may not.
  \end{remark}

  \begin{theorem}[Global log-convexity]\label{thm:logconvex}
  Fix a graph $G$ with $m \ge 1$ edges and define
  $F_G(\theta) = \log LO_G(\theta)$ on $\mathbb{R}^3$. Let
  $\mathbb{E}_{\theta}$ and $\operatorname{Cov}_{\theta}$ denote mean
  and covariance under the edge probability measure
  $w_{uv}(\theta) = e^{\theta^{\top} z_{uv}} / LO_G(\theta)$. Then
  $F_G$ is convex, with
  \begin{equation*}
  \nabla F_G(\theta) = \mathbb{E}_{\theta}[z], \qquad
  \nabla^2 F_G(\theta) = \operatorname{Cov}_{\theta}(z) \succeq 0 .
  \end{equation*}
  In particular
  $\partial_{\gamma} F_G = \mathbb{E}_{\theta}[q] \in [0, 1)$ and
  $\partial^2_{\gamma} F_G = \operatorname{Var}_{\theta}(q) \ge 0$,
  so $\log LO$ is nondecreasing and convex in $\gamma$, strictly
  increasing in $\gamma$ unless every edge of $G$ joins vertices of
  equal degree. Moreover
  $\operatorname{Var}_{\theta}(q) = 0$ at some (equivalently, every)
  $\theta$ if and only if $q_{uv}$ is constant over the edges of $G$:
  this covers regular graphs ($q \equiv 0$), but also, e.g.,
  $(\Delta,\delta)$-biregular graphs, on which every edge carries the
  same nonzero imbalance and $\log LO$ is affine and strictly
  increasing in $\gamma$. $F_G$ is strictly convex if and only if the
  edge feature vectors $\{z_{uv}\}_{uv \in E}$ do not lie in a proper
  affine subspace of $\mathbb{R}^3$; on a regular graph
  $F_G(\theta) = \log m + \theta^{\top} z_0$ is affine and
  $\operatorname{Cov}_{\theta}(z) = 0$.
  \end{theorem}

  \begin{proof}
  Direct differentiation of
  $LO_G(\theta) = \sum_{uv} e^{\theta^{\top} z_{uv}}$ gives
  $\partial_k F_G = \sum_{uv} w_{uv} z_{uv,k} =
  \mathbb{E}_{\theta}[z_k]$ and
  $\partial_k \partial_l F_G =
  \mathbb{E}_{\theta}[z_k z_l] -
  \mathbb{E}_{\theta}[z_k]\,\mathbb{E}_{\theta}[z_l]$, the
  $(k,l)$-entry of $\operatorname{Cov}_{\theta}(z)$, which is positive
  semidefinite; convexity follows. The measure $w(\theta)$ has full
  support on $E$ for every $\theta$, so
  $\operatorname{Cov}_{\theta}(z)$ is singular in a direction $v$ iff
  $v^{\top} z_{uv}$ is constant over all edges, i.e.\ iff the
  $z_{uv}$ lie in a proper affine subspace; this proves the strictness
  criterion. Since $q_{uv} \in [0,1)$,
  $\mathbb{E}_{\theta}[q] \in [0,1)$, and
  $\mathbb{E}_{\theta}[q] = 0$ iff $q_{uv} = 0$ on every edge. On an
  $r$-regular graph all $z_{uv}$ equal
  $z_0 = (2\ln r, \ln 2r, 0)$.
  \end{proof}

  Theorem~\ref{thm:logconvex} places the family in the classical
  exponential-family framework: moving $\gamma$ reweights the edge
  measure toward imbalanced edges at rate
  $\mathbb{E}_{\theta}[q]$, and the Cauchy--Schwarz midpoint
  inequality of Section~\ref{sec:bounds} is precisely midpoint
  log-convexity, i.e.\ a corollary of this theorem.

  \begin{proposition}[Degree-pair information ceiling]\label{prop:collision}
  If two graphs $G$ and $H$ have identical unordered
  edge-degree-pair count vectors, $m_{ij}(G) = m_{ij}(H)$ for all
  $i \le j$, then $LO(G;\theta) = LO(H;\theta)$ for every
  $\theta \in \mathbb{R}^3$, and more generally
  $TI(G) = TI(H)$ for every vertex-degree-based index $TI$. Among the
  $18$ octane isomers exactly $16$ distinct count vectors occur: the
  pairs $\{$3-methylheptane, 4-methylheptane$\}$ and
  $\{$3,4-dimethylhexane, 2-methyl-3-ethylpentane (alphabetized
  IUPAC form: 3-ethyl-2-methylpentane)$\}$ share their
  count vectors while differing in measured properties
  (Table~\ref{tab:collisions}), so every vertex-degree-based
  descriptor, tuned or not, carries a nonzero irreducible error on
  this benchmark.
  \end{proposition}

  \begin{proof}
  By Definition~\ref{def:wuzi}, $LO(G;\theta)$ depends on $G$ only
  through the multiset of unordered edge degree pairs, i.e.\ through
  $(m_{ij}(G))_{i \le j}$; the same holds for every vertex-degree-based
  index~\cite{Rada2026LinearVDB}. The counts for the $18$ octane
  isomers are computed and cross-checked in the reproducibility
  package; both collision pairs have their common count vectors listed
  in Table~\ref{tab:collisions}.
  \end{proof}

  The ceiling is tight in the following sense: distinct count vectors
  are separated by the family at almost every parameter choice.

  \begin{proposition}[Generic discrimination of distinct
  profiles]\label{prop:generic}
  Fix $\Delta \ge 1$ and let $m \neq m'$ be two distinct
  edge-degree-pair count vectors supported on the admissible pairs
  $P_{\Delta}$. Then
  $D(\theta) = \sum_{(i,j) \in P_{\Delta}} (m_{ij} - m'_{ij})\,
  e^{\theta^{\top} z_{ij}}$ is a real-analytic function of $\theta$
  that is not identically zero, and its zero set has Lebesgue measure
  zero in $\mathbb{R}^3$. Consequently, any two graphs with distinct
  count vectors are distinguished by $LO(\cdot\,;\theta)$ for almost
  every $\theta \in \mathbb{R}^3$; only graphs with identical count
  vectors (Proposition~\ref{prop:collision}) are indistinguishable
  for all $\theta$.
  \end{proposition}

  \begin{proof}
  The map $(i,j) \mapsto z_{ij}$ is injective on $P_{\Delta}$, since
  $\ln(ij)$ and $\ln(i+j)$ recover the product and the sum, which
  determine the unordered pair. Choose $v \in \mathbb{R}^3$ such that
  the inner products $v^{\top} z_{ij}$ are pairwise distinct (this
  excludes finitely many hyperplanes of $v$'s). Along the ray
  $\theta = t v$, $D(tv) = \sum_k c_k e^{t \lambda_k}$ with distinct
  exponents $\lambda_k$ and coefficients
  $c_k = m_{i_kj_k} - m'_{i_kj_k}$, not all zero; such an exponential
  sum vanishes identically only if all $c_k = 0$, so $D \not\equiv 0$
  on the ray and hence on $\mathbb{R}^3$. $D$ is a finite sum of
  exponentials of linear functionals, hence real-analytic on
  $\mathbb{R}^3$, and the zero set of a nonzero real-analytic
  function has Lebesgue measure zero.
  \end{proof}

  \section{Bounds for the Loyola family}\label{sec:bounds}
  % =====================================================================

  This section establishes a coherent set of upper and lower bounds
  on $\Wuzi$. The development proceeds from a bracket bound via the
  edge-contribution function $\psi$,  bounds in $(n, m, \Delta, \delta)$
  for the non-negative
  parameter region, Cauchy--Schwarz and
  Jensen-type bounds in the first two Zagreb indices $M_1$ and
  $M_2$,
  a uniform ratio-bound theorem, and bounds via the generalised
  Randi\'c, sum-connectivity, and Sombor indices. Throughout, $G$
  is a simple connected graph with $\delta \le d_v \le \Delta$ for
  every $v \in V(G)$.

  \begin{remark}[Attribution at $\gamma = 0$]\label{rem:gmattr}
  At $\gamma = 0$ the family coincides with $M_{\alpha,\beta}$, for
  which inequalities of all the types developed here are already
  known: $(\Delta,\delta,m)$-bounds with regular/biregular equality
  and two-sided comparisons against $M_1$, $M_2$, $R$, $\chi$, $H$,
  $GA$, $AG$ and $ISI$ appear
  in~\cite{MolinaSigarreta2025GM}, and Jensen-, H\"older- and
  Cauchy--Schwarz-type relations in~\cite{GranadosJMC2025GM,
  PortillaJMC2025GM}. The $\gamma = 0$ specializations of the results
  below are therefore not claimed as new; the content of this section
  is their extension to $\gamma \neq 0$, where the imbalance factor
  $e^{\gamma q_{uv}}$ enters every edge weight, and the tracking of
  how $\gamma$ deforms the equality cases.
  Theorems~\ref{thm:m1_bound}--\ref{thm:m2_alpha} and
  Propositions~\ref{thm:R_chi}--\ref{thm:sombor} are stated at
  $\gamma = 0$ only; they are standard Jensen and Cauchy--Schwarz
  consequences for $M_{\alpha,\beta}$, included to make the
  comparison set self-contained, and are not claimed as new.
  \end{remark}

  The summand in~\eqref{eq:wuzi} depends only on the unordered
  edge degree pair $(d_u, d_v)$. We isolate this dependence as follows:

  \begin{definition}\label{def:psi}
  For real parameters $\alpha, \beta, \gamma$ and positive reals
  $x, y$,
  \begin{equation}\label{eq:psi}
  \edgepsi = (xy)^{\alpha}(x + y)^{\beta}
  \exp\!\left(\gamma\,\tfrac{|x - y|}{x + y}\right).
  \end{equation}
  \end{definition}

  With this notation $\Wuzi = \sum_{uv \in E(G)} \psi(d_u, d_v; \alpha,
  \beta, \gamma)$.

  \begin{proposition}\label{thm:bracket}
  Let $G$ be a simple connected graph of size $m \ge 1$ with
  degrees in $[\delta, \Delta]$. Then
  \begin{equation}\label{eq:bracket}
  m \cdot \psi_{\min}(\delta, \Delta; \alpha, \beta, \gamma)
  \;\le\; \Wuzi \;\le\;
  m \cdot \psi_{\max}(\delta, \Delta; \alpha, \beta, \gamma),
  \end{equation}
  where,
  $\psi_{\min} = \min\{\psi(i, j; \alpha, \beta, \gamma)
                : \delta \le i \le j \le \Delta\}$
  and $\psi_{\max} = \max\{\psi(i, j; \alpha, \beta, \gamma)
                : \delta \le i \le j \le \Delta\}$. The upper
  inequality is an equality if and only if every edge of $G$ has a
  degree pair attaining $\psi_{\max}$, and the lower
  inequality is an equality if and only if every edge has a degree
  pair attaining $\psi_{\min}$.
  \end{proposition}

  \begin{proof}
  Because $G$ has $\delta \le d_v \le \Delta$ for every
  $v \in V(G)$, every edge $uv \in E(G)$ has its degree
  pair $\{d_u, d_v\}$ in the admissible set
  $\mathcal{P}(\delta, \Delta) := \{(i, j) \in \mathbb{Z}^2 : \delta \le i \le j \le
  \Delta\}$.
  By the definitions of $\psi_{\min}$ and $\psi_{\max}$ as the
  minimum and maximum of $\psi$ over $\mathcal{P}(\delta, \Delta)$,
  \[
  \psi_{\min}(\delta, \Delta; \alpha, \beta, \gamma)
  \;\le\;
  \psi(d_u, d_v; \alpha, \beta, \gamma)
  \;\le\;
  \psi_{\max}(\delta, \Delta; \alpha, \beta, \gamma)
  .
  \]
  Summing each of the two per-edge inequalities over the $m$ edges
  of $G$ and using $\sum_{uv \in E(G)} \psi(d_u, d_v; \alpha, \beta,
  \gamma) = \Wuzi$
  gives
  \[
  m \cdot \psi_{\min}(\delta, \Delta; \alpha, \beta, \gamma)
  \;\le\;
  \Wuzi
  \;\le\;
  m \cdot \psi_{\max}(\delta, \Delta; \alpha, \beta, \gamma),
  \]

  For the equality cases, the upper inequality
  $\Wuzi \le m \cdot \psi_{\max}$ is the sum of $m$ per-edge
  inequalities $\psi(d_u, d_v; \alpha, \beta, \gamma) \le \psi_{\max}$;
  this sum is an equality if and only if every per-edge inequality
  is, equivalently every edge $uv$ has a degree pair attaining
  $\psi_{\max}$ over $\mathcal{P}(\delta, \Delta)$. The lower
  equality case is the dual.
  \end{proof}

  \begin{theorem}\label{thm:cs}
  Let $\alpha_1 + \alpha_2 = 2\alpha$, $\beta_1 + \beta_2 = 2\beta$,
  and $\gamma_1 + \gamma_2 = 2\gamma$. Then for every simple
  connected graph $G$ without isolated vertices,
  \[
  LO(G; \alpha, \beta, \gamma)
  \;\le\;
  \sqrt{LO(G; \alpha_1, \beta_1, \gamma_1)
         \cdot LO(G; \alpha_2, \beta_2, \gamma_2)}
  \]
  Equality holds if and only if the ratio
  $\psi(d_u, d_v; \alpha_1, \beta_1, \gamma_1)/\psi(d_u, d_v; \alpha_2,
  \beta_2, \gamma_2)$
  is the same for every edge of $G$.
  \end{theorem}

  \begin{proof}
  For each edge $e = uv \in E(G)$ and each $k \in \{1, 2\}$ write
  \[
  \psi_e^{(k)}
  \;:=\; \psi(d_u, d_v; \alpha_k, \beta_k, \gamma_k)
  \;=\; (d_u d_v)^{\alpha_k}\,(d_u + d_v)^{\beta_k}\,
  \exp\!\left(\gamma_k\, \tfrac{|d_u - d_v|}{d_u + d_v}\right).
  \]
  Both $\psi_e^{(1)}$ and $\psi_e^{(2)}$ are strictly positive on
  every edge of a graph without isolated vertices.

  By the additivity of the exponents and the hypothesis
  $\alpha_1 + \alpha_2 = 2\alpha$,
  $\beta_1 + \beta_2 = 2\beta$, $\gamma_1 + \gamma_2 = 2\gamma$,
  \begin{align*}
  \psi_e^{(1)}\,\psi_e^{(2)}
  &= (d_u d_v)^{\alpha_1 + \alpha_2}\,(d_u + d_v)^{\beta_1 + \beta_2}\,
     \exp\!\left((\gamma_1 + \gamma_2)\, \tfrac{|d_u - d_v|}{d_u +
  d_v}\right) \\
  &= (d_u d_v)^{2\alpha}\,(d_u + d_v)^{2\beta}\,
     \exp\!\left(2\gamma\, \tfrac{|d_u - d_v|}{d_u + d_v}\right) \\
  &= \psi(d_u, d_v; \alpha, \beta, \gamma)^{2}.
  \end{align*}
  Taking square roots gives
  $\sqrt{\psi_e^{(1)}\,\psi_e^{(2)}} = \psi(d_u, d_v; \alpha, \beta,
  \gamma)$
  for every edge. Apply the Cauchy--Schwarz inequality to the
  positive sequences $a_e = \sqrt{\psi_e^{(1)}}$,
  $b_e = \sqrt{\psi_e^{(2)}}$ indexed over $E(G)$:
  \[
  \sum_{e \in E(G)} a_e\, b_e
  \;\le\;
  \sqrt{\sum_{e \in E(G)} a_e^{2}}\,
  \sqrt{\sum_{e \in E(G)} b_e^{2}}.
  \]
  Substituting $a_e b_e = \sqrt{\psi_e^{(1)} \psi_e^{(2)}}
     = \psi(d_u, d_v; \alpha, \beta, \gamma)$,
  $\sum_e a_e^2 = \sum_e \psi_e^{(1)} = LO(G; \alpha_1, \beta_1,
  \gamma_1)$,
  and $\sum_e b_e^2 = \sum_e \psi_e^{(2)} = LO(G; \alpha_2, \beta_2,
  \gamma_2)$
  gives
  \begin{multline*}
  \Wuzi
  \;=\; \sum_{e \in E(G)} \psi(d_u, d_v; \alpha, \beta, \gamma)\\
  \;\le\;
  \sqrt{LO(G; \alpha_1, \beta_1, \gamma_1)\,
        LO(G; \alpha_2, \beta_2, \gamma_2)},
  \end{multline*}
  which is the inequality.

  The Cauchy--Schwarz inequality is an equality if and only if the
  sequences $(a_e)$ and $(b_e)$ are proportional, i.e.\ there exists
  a constant $\lambda \ge 0$ such that $b_e = \lambda\, a_e$ for
  every $e \in E(G)$, equivalently
  $\psi_e^{(2)}/\psi_e^{(1)} = \lambda^{2}$ is the same constant on
  every edge.
  \end{proof}

  \begin{proposition}\label{thm:upper_Delta}
  Let $G$ be a connected graph of size $m$ and maximum degree $\Delta$. For
  $\alpha, \beta \ge 0$ and $\gamma \ge 0$,
  \begin{equation}\label{eq:upper_Delta}
  \Wuzi \;\le\; m \cdot \Delta^{2\alpha} \cdot (2\Delta)^{\beta}
  \cdot \exp\!\left(\gamma\,\frac{\Delta - 1}{\Delta + 1}\right).
  \end{equation}
  \end{proposition}

  \begin{proof}
  Fix an edge $uv \in E(G)$. By the definition of $\Delta$ we have
  $1 \le d_u, d_v \le \Delta$. We bound the three multiplicative
  factors of $\psi(d_u, d_v; \alpha, \beta, \gamma)$ separately.

  \emph{First factor.} Because $d_u, d_v \le \Delta$ we obtain
  $d_u d_v \le \Delta^{2}$. Since $\alpha \ge 0$ and the map
  $t \mapsto t^{\alpha}$ is non-decreasing on $\mathbb{R}_{> 0}$,
  $(d_u d_v)^{\alpha} \le \Delta^{2\alpha}$.

  \emph{Second factor.} Similarly $d_u + d_v \le 2\Delta$, and the
  map $t \mapsto t^{\beta}$ is non-decreasing for $\beta \ge 0$, so
  $(d_u + d_v)^{\beta} \le (2\Delta)^{\beta}$.

  \emph{Third factor.} The normalised imbalance
  $\mathrm{imb}(x, y) = |x - y|/(x + y)$ is maximised over the
  admissible set $\{(i, j) : 1 \le i \le j \le \Delta\}$ at the
  pair $(1, \Delta)$, where
  $\mathrm{imb}(1, \Delta) = (\Delta - 1)/(\Delta + 1)$. Hence
  $|d_u - d_v|/(d_u + d_v) \le (\Delta - 1)/(\Delta + 1)$ for every
  edge. Because $\gamma \ge 0$ and $\exp$ is non-decreasing,
  $\exp(\gamma\,|d_u - d_v|/(d_u + d_v)) \le \exp(\gamma(\Delta - 1)/(\Delta + 1))$.
  Multiplying the three factor-wise bounds and summing over the
  $m$ edges of $G$ yields~\eqref{eq:upper_Delta}. For $\gamma > 0$
  and $\Delta \ge 2$ the bound is attained if and only if
  $\alpha = \beta = 0$ and every edge has degree pair $(1, \Delta)$
  (for connected $G$, $G \cong K_{1,\Delta}$); if $\alpha > 0$ or
  $\beta > 0$ the three factor-wise maxima are not simultaneously
  attainable and the inequality is strict.
  \end{proof}

  \begin{proposition}\label{thm:lower_delta}
  Let $G$ be a connected graph of size $m$ and minimum degree $\delta$. For
  $\alpha, \beta \ge 0$ and $\gamma \ge 0$,
  \begin{equation}\label{eq:lower_delta}
  \Wuzi \;\ge\; m \cdot \delta^{2\alpha} \cdot (2\delta)^{\beta},
  \end{equation}
  with equality if and only if $G$ is $\delta$-regular, except in the
  degenerate case $(\alpha, \beta, \gamma) = (0, 0, 0)$, where
  $\Wuzi \equiv m$ and the bound holds with equality for every graph.
  \end{proposition}

  \begin{proof}
  Fix an arbitrary edge $uv \in E(G)$. By the definition of
  $\delta = \delta(G)$ we have $d_u, d_v \ge \delta$, and we bound
  the three multiplicative factors of
  $\psi(d_u, d_v; \alpha, \beta, \gamma)$ separately. Because
  $d_u, d_v \ge \delta$ and $\alpha \ge 0$, $(d_u d_v)^{\alpha} \ge
  \delta^{2\alpha}$; similarly $(d_u + d_v)^{\beta} \ge (2\delta)^{\beta}$;
  and since $\gamma \ge 0$ and the imbalance is non-negative,
  $\exp(\gamma\,|d_u - d_v|/(d_u + d_v)) \ge 1$. Multiplying the three
  lower bounds gives $\psi(d_u, d_v; \alpha, \beta, \gamma) \ge
  \delta^{2\alpha}(2\delta)^{\beta}$ for every edge, and summing over
  the $m$ edges yields~\eqref{eq:lower_delta}.

  For the equality case, if $(\alpha, \beta, \gamma) = (0, 0, 0)$ then
  $LO(G; 0, 0, 0) = m$ and the lower bound is also $m$, so equality
  holds for every graph. Assume now that
  $(\alpha, \beta, \gamma) \neq (0, 0, 0)$. Equality in the summed
  inequality requires equality in every active per-edge factor. If
  $\alpha > 0$ or $\beta > 0$, this forces $d_u = d_v = \delta$ on every
  edge. If $\alpha = \beta = 0$ and $\gamma > 0$, equality in the
  exponential factor requires $d_u = d_v$ on every edge; since $G$ is
  connected, this implies that $G$ is regular, hence $\delta$-regular.
  Conversely, if $G$ is $\delta$-regular then every edge has degree pair
  $(\delta, \delta)$ and equality holds. Thus, outside the degenerate
  case $(0, 0, 0)$, equality holds if and only if $G$ is
  $\delta$-regular.
  \end{proof}

  \begin{theorem}\label{thm:m1_bound}
  Let $\alpha = \gamma = 0$ and $\beta \ge 1$. Then
  \[
  LO(G; 0, \beta, 0)
  \;=\; \sum_e (d_u + d_v)^{\beta}
  \;\ge\; m^{1 - \beta}\, M_1(G)^{\beta},
  \]
  where $M_1(G) = \sum_{v \in V} d_v^2 = \sum_{uv \in E}(d_u + d_v)$.
  At $\beta = 1$ both sides equal $M_1(G)$ identically; for
  $\beta > 1$ equality holds if and only if $d_u + d_v$ is constant
  on every edge of $G$.
  \end{theorem}

  \begin{proof}
  Setting $\alpha = \gamma = 0$ in Definition~\ref{def:wuzi}
  collapses the edge-contribution function to
  $\psi(d_u, d_v; 0, \beta, 0) = (d_u + d_v)^{\beta}$, so
  $LO(G; 0, \beta, 0) = \sum_{uv \in E(G)} (d_u + d_v)^{\beta}$.
  The function $\varphi(t) = t^{\beta}$ satisfies
  $\varphi''(t) = \beta(\beta - 1)\, t^{\beta - 2} \ge 0$ for
  $\beta \ge 1$; hence $\varphi$ is convex on
  $\mathbb{R}_{> 0}$, with strict convexity for $\beta > 1$. Apply
  Jensen's inequality (convex case) to the $m = |E(G)|$ positive
  values $\{d_u + d_v : uv \in E(G)\}$ with uniform weights $1/m$:
  \[
  \left(\frac{M_1(G)}{m}\right)^{\!\beta}
  = \varphi\!\left(\frac{1}{m}\sum_{uv} (d_u + d_v)\right)
  \;\le\;
  \frac{1}{m}\sum_{uv} (d_u + d_v)^{\beta}
  = \frac{LO(G; 0, \beta, 0)}{m},
  \]
  using $\sum_{uv \in E(G)} (d_u + d_v) = M_1(G)$. Multiplying both
  sides by $m > 0$ yields the bound. At $\beta = 1$, $\varphi$ is
  affine and Jensen is the identity $LO(G; 0, 1, 0) = M_1(G)$; for
  $\beta > 1$ strict convexity gives equality if and only if
  $d_u + d_v$ is constant over every edge.
  \end{proof}

  \begin{theorem}\label{thm:m1_concave}
  Let $\alpha = \gamma = 0$ and $0 < \beta < 1$. Then
  $LO(G; 0, \beta, 0) = \sum_e (d_u + d_v)^{\beta} \le m^{1 - \beta}\, M_1(G)^{\beta}$,
  with equality if and only if $d_u + d_v$ is constant on every
  edge of $G$.
  \end{theorem}

  \begin{proof}
  The function $t \mapsto t^{\beta}$ is strictly concave on
  $\mathbb{R}_{> 0}$ for $\beta \in (0, 1)$. Jensen's inequality in
  the concave case applied to the values $\{d_u + d_v : uv \in E(G)\}$
  gives $\frac{1}{m} \sum_e (d_u + d_v)^{\beta} \le ( M_1(G) / m )^{\beta}$,
  using $M_1(G) = \sum_{uv} (d_u + d_v)$. Multiplying by $m$ gives
  the upper bound; strict concavity gives equality if and only if
  $d_u + d_v$ is constant on every edge.
  \end{proof}

  \begin{theorem}\label{thm:m2_alpha}
  Let $\beta = \gamma = 0$ and $\alpha \ge 1$. Then
  \[
  LO(G; \alpha, 0, 0)
  \;=\; \sum_{uv \in E(G)} (d_u d_v)^{\alpha}
  \;=\; R_{\alpha}(G)
  \;\ge\; m^{\,1 - \alpha}\, M_2(G)^{\alpha},
  \]
  where $M_2(G) = \sum_{uv \in E(G)} d_u d_v$ is the second Zagreb
  index and $R_{\alpha}(G) = \sum_{uv \in E(G)} (d_u d_v)^{\alpha}$
  is the general Randi\'c index. At $\alpha = 1$ both sides equal
  $M_2(G)$ identically; for $\alpha > 1$ equality holds if and only
  if the product $d_u d_v$ is constant on every edge of $G$.
  \end{theorem}

  \begin{proof}
  Setting $\beta = \gamma = 0$ collapses the edge contribution to
  $(d_u d_v)^{\alpha}$, so $LO(G; \alpha, 0, 0) = R_{\alpha}(G)$. The function
  $\varphi(t) = t^{\alpha}$ has
  $\varphi''(t) = \alpha(\alpha - 1)\, t^{\alpha - 2} \ge 0$ for
  $\alpha \ge 1$ and is convex, strictly for $\alpha > 1$. Jensen
  (convex case) on $\{d_u d_v : uv \in E(G)\}$ with weights $1/m$
  gives $(M_2(G)/m)^{\alpha} \le R_{\alpha}(G)/m$, using
  $\sum_{uv} d_u d_v = M_2(G)$. Multiplying by $m$ yields the bound;
  $\alpha = 1$ is the affine identity, and for $\alpha > 1$ strict
  convexity gives equality if and only if $d_u d_v$ is constant on
  every edge.
  \end{proof}

  \begin{proposition}\label{thm:R_chi}
  Let $\gamma = 0$, and define
  $R_a(G) = \sum_{uv \in E(G)} (d_u d_v)^{a}$ and
  $\chi_b(G) = \sum_{uv \in E(G)} (d_u + d_v)^{b}$ (so $R = R_{-1/2}$
  and $\chi = \chi_{-1/2}$). Then for every graph $G$,
  \begin{equation}\label{eq:R_chi_bound}
  LO(G; \alpha, \beta, 0)
  \;\le\; \sqrt{R_{2\alpha}(G) \cdot \chi_{2\beta}(G)}.
  \end{equation}
  In particular, at $(\alpha, \beta) = (-1/4, -1/4)$ this
  specialises to
  \[
  LO(G; -1/4, -1/4, 0) \le \sqrt{R(G)\,\chi(G)}.
  \]
  \end{proposition}

  \begin{proof}
  Apply Theorem~\ref{thm:cs} with
  $(\alpha_1, \beta_1, \gamma_1) = (2\alpha, 0, 0)$ and
  $(\alpha_2, \beta_2, \gamma_2) = (0, 2\beta, 0)$, whose components
  sum to $(2\alpha, 2\beta, 0)$, meeting the hypotheses since
  $\gamma = 0$. By Definition~\ref{def:wuzi},
  $LO(G; 2\alpha, 0, 0) = R_{2\alpha}(G)$ and
  $LO(G; 0, 2\beta, 0) = \chi_{2\beta}(G)$, so Theorem~\ref{thm:cs}
  gives~\eqref{eq:R_chi_bound}. The equality case is that of
  Theorem~\ref{thm:cs}: the ratio
  $(d_u d_v)^{2\alpha}/(d_u + d_v)^{2\beta}$ is constant over every
  edge.
  \end{proof}

  \begin{proposition}\label{thm:sombor}
  For $\alpha = 0$, $0 < \beta < 1$, and $\gamma = 0$, the Loyola
  quantity $LO(G; 0, \beta, 0) = \chi_{\beta}(G)$ satisfies
  \begin{equation}\label{eq:sombor_bound}
  LO(G; 0, \beta, 0)
  \;\le\; m^{1 - \beta} \cdot \bigl(\sqrt{2}\, SO(G)\bigr)^{\beta},
  \end{equation}
  where $SO(G) = \sum_{uv \in E(G)} \sqrt{d_u^2 + d_v^2}$ is the
  Sombor index.
  \end{proposition}

  \begin{proof}
  For positive reals $x, y$, $x^2 + y^2 \ge (x + y)^2/2$ follows from
  $(x + y)^2 \le 2(x^2 + y^2)$ by the AM--GM inequality
  $xy \le (x^2 + y^2)/2$. Taking square roots with $x = d_u, y = d_v$
  and summing over $E(G)$ gives
  $SO(G) \ge M_1(G)/\sqrt{2}$, i.e.\ $M_1(G) \le \sqrt{2}\, SO(G)$.
  For $0 < \beta < 1$, the concave Jensen bound
  (Theorem~\ref{thm:m1_concave}) gives
  $LO(G; 0, \beta, 0) \le m^{1 - \beta}\, M_1(G)^{\beta}$, and since
  $t \mapsto t^{\beta}$ is non-decreasing for $\beta > 0$,
  substituting $M_1(G) \le \sqrt{2}\, SO(G)$
  yields~\eqref{eq:sombor_bound}.
  \end{proof}

  The next theorem gives a single uniform template for bounding
  $\Wuzi$ both above and below in terms of any classical
  bond-incident-degree index whose edge contribution is strictly
  positive on the admissible degree pairs of the graph.

  \begin{theorem}\label{thm:ratio_bound}
  Let $G$ be a simple graph with maximum degree at most $\Delta$, and let
  $I_h(G) = \sum_{uv \in E(G)} h(d_u, d_v)$ be a BID index with
  edge contribution $h(i, j)$ strictly positive on every degree pair
  $(i, j)$ that appears in $G$. Define
  \[
  c_{\min}^{h} = \min_{(i, j)}\,\frac{\psi(i, j; \alpha, \beta,
  \gamma)}{h(i, j)},
  \qquad
  c_{\max}^{h} = \max_{(i, j)}\,\frac{\psi(i, j; \alpha, \beta,
  \gamma)}{h(i, j)},
  \]
  where the extrema range over the degree pairs actually appearing in
  $G$.
  Then
  \begin{equation}\label{eq:ratio_bound}
  c_{\min}^{h} \cdot I_h(G)
  \;\le\; \Wuzi
  \;\le\; c_{\max}^{h} \cdot I_h(G),
  \end{equation}
  with equality on either side if and only if every edge of $G$
  attains the corresponding extremum of $\psi/h$.
  \end{theorem}

  \begin{proof}
  Fix an edge $uv \in E(G)$ and define
  \[
  r_{uv} := \psi(d_u, d_v; \alpha, \beta, \gamma)/h(d_u, d_v) > 0.
  \]
  By definition $c_{\min}^{h} \le r_{uv} \le c_{\max}^{h}$, so
  multiplying by $h(d_u, d_v) > 0$ gives the per-edge bounds
  $c_{\min}^{h} h(d_u, d_v) \le \psi(d_u, d_v; \alpha, \beta, \gamma)$
  and $\psi(d_u, d_v; \alpha, \beta, \gamma)
  \le c_{\max}^{h} h(d_u, d_v)$. Summing over $E(G)$ and using
  $\sum_{uv} h(d_u, d_v) = I_h(G)$ and $\sum_{uv} \psi = \Wuzi$
  gives~\eqref{eq:ratio_bound}; the equality conditions are that
  every edge attains the corresponding extremum of $\psi/h$ over the
  appearing pairs.
  \end{proof}

  Theorem~\ref{thm:ratio_bound} is a \emph{sharp graph-specific}
  comparison: its constants depend on the degree pairs that actually
  appear in $G$. A uniform version follows immediately by enlarging
  the range of the extrema.

  \begin{corollary}[Uniform $\Delta$-bounded version]\label{cor:ratio_uniform}
  With the hypotheses of Theorem~\ref{thm:ratio_bound}, but with $h$
  strictly positive on all of
  $P_{\Delta} = \{(i,j) : 1 \le i \le j \le \Delta\}$ and the extrema
  taken over $P_{\Delta}$,
  inequality~\eqref{eq:ratio_bound} holds simultaneously for every
  graph of maximum degree at most $\Delta$, with constants independent
  of the individual graph. Equality on either side holds if and only
  if every edge of $G$ realises a degree pair attaining the
  corresponding extremum over $P_{\Delta}$; the uniform constants are
  in general strictly worse than the graph-specific ones, since the
  extrema are taken over a superset.
  \end{corollary}


  \begin{proposition}\label{prop:gamma_eq_0_or_inf}
  At $\gamma = 0$, $LO(G; \alpha, \beta, 0)$ coincides with the
  two-parameter $(\alpha, \beta)$ degree-based index
  $\sum_{uv \in E(G)} (d_u d_v)^{\alpha}(d_u + d_v)^{\beta}$.
  Let $E_{\max}(G)$ denote the set of edges realising
  $\mathrm{imb}_{\max}(G) = \max_{uv} \mathrm{imb}(d_u, d_v)$
  and $E_{\min}(G)$ the set realising
  $\mathrm{imb}_{\min}(G) = \min_{uv} \mathrm{imb}(d_u, d_v)$.
  Then as $\gamma \to +\infty$,
  \[
  \frac{\Wuzi}{e^{\gamma\, \mathrm{imb}_{\max}(G)}}
  \;\longrightarrow\;
  \sum_{uv \in E_{\max}(G)} (d_u d_v)^{\alpha}(d_u + d_v)^{\beta},
  \]
  and as $\gamma \to -\infty$,
  $\Wuzi / e^{\gamma\, \mathrm{imb}_{\min}(G)} \longrightarrow
  \sum_{uv \in E_{\min}(G)} (d_u d_v)^{\alpha}(d_u + d_v)^{\beta}$.
  \end{proposition}

  \begin{proof}
  At $\gamma = 0$ the exponential factor is $1$ on every edge.
  For the asymptotics, write $\mu(G) := \mathrm{imb}_{\max}(G)$ and
  $\varepsilon_{uv} := \mu(G) - \mathrm{imb}(d_u, d_v) \ge 0$, so that
  $\Wuzi / e^{\gamma \mu(G)} = \sum_{uv} (d_u d_v)^{\alpha}(d_u +
  d_v)^{\beta} e^{-\gamma \varepsilon_{uv}}$. For $uv \in E_{\max}(G)$,
  $\varepsilon_{uv} = 0$ and the summand is $\gamma$-independent; for
  $uv \notin E_{\max}(G)$, $\varepsilon_{uv} > 0$ and $e^{-\gamma \varepsilon_{uv}}
  \to 0$ as $\gamma \to +\infty$. Since $E(G)$ is finite the limit
  is termwise, giving the identity. The $\gamma \to -\infty$ case is
  dual.
  \end{proof}




% =====================================================================
\section{Sensitivity analysis}\label{sec:sensitivity}
% =====================================================================

We benchmark the Loyola family against the classical indices it
reduces to, on three standard testbeds, following the chemometric
benchmark of the diminished Sombor
study~\cite{MovahediGutmanRedzepovicFurtula2026}: the $18$ octane
isomers (five NIST physicochemical properties), the $106$
non-isomorphic trees of order $10$, and the $75$ decane isomers. The comparison set
is restricted to the exact reductions of Section~\ref{sec:intro} ---
$M_1, M_2, HM, {}^{m}\!M_2, R, \chi, H/2, ISI, GA, AG$ --- together
with the two fixed pure-$\gamma$ benchmark points $LO(0, 0, 1)$
and $LO(0, 0, 2)$ chosen for this study. No parameter in this comparison
is tuned against the property data --- the effect of data-driven
tuning is examined separately, as a negative control, in
Section~\ref{sec:tuning}; the edge count $m = LO(0, 0, 0)$ is omitted
as it is constant on the octanes. All comparisons in this section are
descriptive; given the number of properties, indices, and tuning
candidates examined, no confirmatory inference is claimed.

\subsection{Property correlations on the octane isomers}\label{sec:octanes}

Table~\ref{tab:lo_octane} reports the Pearson correlation $r$ between
each index and each of the five physicochemical properties of the
$18$ octane isomers (boiling point $T_B$, enthalpy of formation
$\Delta H_f$, enthalpy of vaporization $\Delta H_{\mathrm{vap}}$,
entropy $S$, acentric factor $\omega$; property data from the NIST
WebBook~\cite{NISTWebBook}, compiled for internal consistency and comparison with
diminished-Sombor-style benchmarks, accessed 2024--2025). Units are: $T_B$ in
$^{\circ}$C; $\Delta H_f$ and $\Delta H_{\mathrm{vap}}$ in
kcal\,mol$^{-1}$ (the latter at $298$~K); $S$ in
cal\,mol$^{-1}$K$^{-1}$; and $\omega$ dimensionless. Values for
$\Delta H_{\mathrm{vap}}$ and entropy $S$ vary slightly across
compilations; all comparisons here are made within this single,
internally consistent dataset. Calculations use this author-compiled
property table attributed to the NIST WebBook. All $90$
molecule-level values were traced against authoritative sources
(NIST WebBook first; KDB-derived compilations for the acentric
factor, which NIST does not list), and classified under a strict
taxonomy in which ``verified'' requires agreement at the value's
reported precision: $3$ match exactly, $34$ match at reported
precision after unit conversion (gas phase confirmed for formation
enthalpies and entropies), $32$ agree within the audit tolerances or
the source's stated uncertainty without matching at reported
precision, $13$ lie within explained variation among authoritative
determinations, $7$ could not be verified against a reachable
authoritative value (six gas-phase entropies absent from
the NIST WebBook, and the $298$\,K vaporization enthalpy of
2,2,3,3-tetramethylbutane, which is solid at room temperature so
that sublimation rather than vaporization data are tabulated), and
one value is a marginal conflict: the acentric factors used here for
2,2,4-trimethylpentane ($0.305$ vs.\ compiled $0.303$) and
2,2,3,3-tetramethylbutane ($0.247$ vs.\ compiled $0.251$) differ
from the KDB compilation by up to $0.004$. Substituting the
alternative compiled acentric factors changes several of the
$\omega$ correlations in Table~\ref{tab:lo_octane} at the third
decimal (e.g.\ $M_2$ moves from $-0.988$ to $-0.986$) but does not
change the property-wise winning descriptor or any substantive
interpretation; the full sensitivity computation and the
machine-readable provenance trace accompany the reproducibility
package. The benchmark values themselves are preserved as supplied
for comparability.

\begin{table}[H]
\centering
\caption{Pearson correlation $r$ between the Loyola reductions, the fixed pure-$\gamma$ benchmark points and the five physicochemical properties of the $18$ octane isomers. Bold flags the best $|r|$ per column.}
\label{tab:lo_octane}
\small
\begin{tabular}{lrrrrr}
\toprule
Index & $T_B$ & $\Delta H_f$ & $\Delta H_{\mathrm{vap}}$ & $S$ & $\omega$\\
\midrule
$M_1 = LO(0,1,0)$ & $-0.718$ & $-0.762$ & $-0.912$ & $-0.973$ & $-0.970$\\
$M_2 = LO(1,0,0)$ & $-0.497$ & $-0.542$ & $-0.772$ & $-0.940$ & $\mathbf{-0.988}$\\
$HM = LO(0,2,0)$ & $-0.654$ & $-0.710$ & $-0.873$ & $\mathbf{-0.975}$ & $-0.982$\\
${}^{m}\!M_2 = LO(-1,0,0)$ & $\mathbf{+0.855}$ & $\mathbf{+0.891}$ & $+0.937$ & $+0.863$ & $+0.794$\\
$R = LO(-1/2,0,0)$ & $+0.819$ & $+0.850$ & $+0.954$ & $+0.933$ & $+0.898$\\
$\chi = LO(0,-1/2,0)$ & $+0.800$ & $+0.832$ & $+0.951$ & $+0.949$ & $+0.924$\\
$H/2 = LO(0,-1,0)$ & $+0.821$ & $+0.849$ & $+0.956$ & $+0.933$ & $+0.900$\\
$ISI = LO(1,-1,0)$ & $-0.005$ & $-0.000$ & $-0.331$ & $-0.601$ & $-0.738$\\
$GA = 2\,LO(1/2,-1,0)$ & $+0.822$ & $+0.858$ & $+0.957$ & $+0.941$ & $+0.906$\\
$AG = \tfrac12 LO(-1/2,1,0)$ & $-0.816$ & $-0.859$ & $-0.950$ & $-0.939$ & $-0.900$\\
$LO(0,0,1)$ & $-0.828$ & $-0.832$ & $-0.966$ & $-0.931$ & $-0.929$\\
$LO(0,0,2)$ & $-0.829$ & $-0.846$ & $\mathbf{-0.967}$ & $-0.939$ & $-0.925$\\
\bottomrule
\end{tabular}
\end{table}

The best correlation per property is split among the reductions and
the pure-$\gamma$ points: the modified second Zagreb index
$LO(-1, 0, 0)$ is strongest on $T_B$ and $\Delta H_f$, the
hyper-Zagreb index $LO(0, 2, 0)$ on $S$, the second Zagreb index
$LO(1, 0, 0)$ on $\omega$, and the fixed pure-$\gamma$ benchmark point
$LO(0, 0, 2)$ on $\Delta H_{\mathrm{vap}}$ ($|r| = 0.967$). At the
exact-reduction triples the Loyola correlations coincide with the
corresponding classical indices to floating-point precision, as they
must. For each property, a descriptive post-selection paired bootstrap
interval for the observed top-two $|r|$ difference ($B = 10\,000$,
resampling the $18$ isomers jointly) included zero; these intervals
were not multiplicity-adjusted across the five properties and do not
establish equality or equivalence. Individual bootstrap half-widths
range from about $0.02$ to $0.10$ across the headline cells. The value of $LO(0, 0, 2)$ on $\Delta H_{\mathrm{vap}}$
is a fixed benchmark point, not a tuned result. Recomputing each bolded correlation with single isomers deleted
shifts $|r|$ by at most $0.027$, indicating that no single isomer
drives the leading correlations.

Figure~\ref{fig:octanes} shows the signed-$r$ heatmap of the
compared indices against the five properties.

\begin{figure}[H]
\centering
\safefig{fig_lo_prediction_correlation_heatmap_v3.pdf}{0.9}
\caption{Signed Pearson $r$ between the Loyola reductions, the fixed pure-$\gamma$ benchmark points and the five octane physicochemical properties on the $18$ isomers.}
\label{fig:octanes}
\end{figure}

On the octanes the pure-$\gamma$ points are highly collinear with
the classical reductions (maximum $|r|$ with a reduction equal to
$0.982$ for $LO(0, 0, 1)$ and $0.993$ for $LO(0, 0, 2)$).

\paragraph{Descriptor redundancy.} The redundancy is in fact
systemic, and it caps how much meaning any small $|r|$ difference in
Table~\ref{tab:lo_octane} can carry. Across the twelve compared
descriptors on the $18$ isomers, fourteen descriptor pairs have
$|r| \ge 0.99$ (the extreme being $R$ versus $H/2$ at $r = 0.9997$),
and a principal component analysis of the standardized $18 \times 12$
descriptor matrix places $90.3\%$ of the variance on the first
component and $99.05\%$ on the first two, with an effective rank
(participation ratio of the eigenvalue spectrum) of $1.2$. The
twelve-index comparison set is thus essentially a two-dimensional
descriptor space, and differences in the third decimal of $|r|$
between near-collinear descriptors should not be read as meaningful
rankings.

\paragraph{The information ceiling realised on this benchmark.}
Proposition~\ref{prop:collision} is not hypothetical on these data:
the $18$ isomers realise only $16$ distinct edge-degree-pair count
vectors, and the two collision pairs of Table~\ref{tab:collisions}
are indistinguishable by \emph{every} vertex-degree-based index at
\emph{every} parameter choice, while their measured properties
differ (by up to $1.2\,^{\circ}$C in $T_B$ and $1.94$ units in $S$
within the first pair, and $2.1\,^{\circ}$C in $T_B$ within the
second). Every correlation in Table~\ref{tab:lo_octane}, and every
tuned variant in the sections below, operates above this
irreducible floor.

\begin{table}[H]
\centering
\caption{The two degree-pair collision pairs among the $18$ octane
isomers. Each pair shares its complete edge-degree-pair count vector
(shown as multiplicities of the unordered pairs), hence is
indistinguishable by every vertex-degree-based index for all
parameter values, yet differs in the measured properties.}
\label{tab:collisions}
\small
\begin{tabular}{lrrrrr}
\toprule
Isomer & $T_B$ & $\Delta H_f$ & $\Delta H_{\mathrm{vap}}$ & $S$ & $\omega$\\
\midrule
\multicolumn{6}{l}{\emph{shared count vector} $2(1,2),\,1(1,3),\,2(2,2),\,2(2,3)$:}\\
3-methylheptane & $118.9$ & $-50.82$ & $9.52$ & $111.26$ & $0.371$\\
4-methylheptane & $117.7$ & $-50.69$ & $9.48$ & $109.32$ & $0.372$\\
\midrule
\multicolumn{6}{l}{\emph{shared count vector} $2(1,2),\,2(1,3),\,2(2,3),\,1(3,3)$:}\\
3,4-dimethylhexane & $117.7$ & $-50.91$ & $9.32$ & $106.59$ & $0.340$\\
2-methyl-3-ethylpentane & $115.6$ & $-50.48$ & $9.21$ & $106.06$ & $0.330$\\
\bottomrule
\end{tabular}
\end{table}

\subsection{Effect of parameter tuning near the best indices}\label{sec:tuning}

The benchmark of Section~\ref{sec:octanes} tunes no parameters: it
compares only the exact reductions and the two fixed pure-$\gamma$
benchmark points. It is nonetheless natural to ask whether small,
data-driven adjustments of $(\alpha, \beta, \gamma)$ near the
best-correlating indices could improve the fit. We probe this directly
as a negative control. For each of the four anchors carrying the
strongest column correlations in Table~\ref{tab:lo_octane} ---
$M_2 = LO(1, 0, 0)$, $HM = LO(0, 2, 0)$, ${}^{m}\!M_2 = LO(-1, 0, 0)$,
and the pure-$\gamma$ point $LO(0, 0, 1)$ --- we draw $50$ random
parameter triples within $\pm 0.3$ (to two decimal places) of the
anchor and correlate each against the five octane properties.
Table~\ref{tab:lo_tuning} reports, for every anchor--property cell, the
anchor correlation $|r|_0$, the best of the $50$ perturbations in
sample $|r|_{50}$, the in-sample gain $\Delta_{\mathrm{in}}$, the
percentile-bootstrap half-width ($B = 10\,000$), the leave-one-out
correlation of the anchor $\mathrm{LOO}_0$, and a \emph{nested}
leave-one-out correlation $\mathrm{LOO}_{50}$ in which, within each of
the four globally selected anchor pools, the best of the $50$
candidates is selected using the $17$ training isomers of each fold
only and evaluated on the held-out one, with the resulting
out-of-sample gain $\Delta_{\mathrm{out}}$. The anchors themselves
were chosen from the complete-data correlations of
Table~\ref{tab:lo_octane}, and the $18$ pooled fold predictions share
heavily overlapping training sets, so these values are descriptive
rather than independent replications.

\begin{table}[H]
\centering
\caption{Decimal parameter tuning near the four best-correlating anchors on the $18$ octane isomers ($50$ random triples within $\pm 0.3$ of each anchor). Columns: anchor correlation $|r|_0$; best-of-$50$ in sample $|r|_{50}$; in-sample gain $\Delta_{\mathrm{in}}$; percentile-bootstrap half-width $\pm$boot ($B = 10\,000$); leave-one-out correlation of the anchor $\mathrm{LOO}_0$; foldwise nested leave-one-out correlation of the tuned family, computed within the globally selected anchor pool, $\mathrm{LOO}_{50}$; out-of-sample gain $\Delta_{\mathrm{out}} = \mathrm{LOO}_{50} - \mathrm{LOO}_0$.}
\label{tab:lo_tuning}
\scriptsize
\setlength{\tabcolsep}{4pt}
\begin{tabular}{lrrrrrrr}
\toprule
Property & $|r|_0$ & $|r|_{50}$ & $\Delta_{\mathrm{in}}$ & $\pm$boot & $\mathrm{LOO}_0$ & $\mathrm{LOO}_{50}$ & $\Delta_{\mathrm{out}}$\\
\midrule
\multicolumn{8}{l}{\emph{near} $M_2 = LO(1,0,0)$}\\
$T_B$ & $0.497$ & $0.598$ & $+0.100$ & $0.353$ & $0.331$ & $0.458$ & $+0.127$\\
$\Delta H_f$ & $0.542$ & $0.635$ & $+0.093$ & $0.358$ & $0.395$ & $0.540$ & $+0.145$\\
$\Delta H_{\mathrm{vap}}$ & $0.772$ & $0.839$ & $+0.068$ & $0.231$ & $0.714$ & $0.786$ & $+0.073$\\
$S$ & $0.940$ & $0.960$ & $+0.021$ & $0.067$ & $0.919$ & $0.943$ & $+0.024$\\
$\omega$ & $0.988$ & $0.995$ & $+0.007$ & $0.018$ & $0.985$ & $0.985$ & $\phantom{+}0.000$\\
\midrule
\multicolumn{8}{l}{\emph{near} $HM = LO(0,2,0)$}\\
$T_B$ & $0.654$ & $0.740$ & $+0.085$ & $0.257$ & $0.545$ & $0.647$ & $+0.102$\\
$\Delta H_f$ & $0.710$ & $0.786$ & $+0.076$ & $0.255$ & $0.645$ & $0.741$ & $+0.097$\\
$\Delta H_{\mathrm{vap}}$ & $0.873$ & $0.920$ & $+0.047$ & $0.122$ & $0.837$ & $0.894$ & $+0.057$\\
$S$ & $0.975$ & $0.976$ & $+0.002$ & $0.027$ & $0.967$ & $0.960$ & $-0.007$\\
$\omega$ & $0.982$ & $0.989$ & $+0.007$ & $0.025$ & $0.979$ & $0.982$ & $+0.003$\\
\midrule
\multicolumn{8}{l}{\emph{near} ${}^{m}\!M_2 = LO(-1,0,0)$}\\
$T_B$ & $0.855$ & $0.864$ & $+0.009$ & $0.096$ & $0.781$ & $0.759$ & $-0.021$\\
$\Delta H_f$ & $0.891$ & $0.897$ & $+0.006$ & $0.087$ & $0.865$ & $0.837$ & $-0.028$\\
$\Delta H_{\mathrm{vap}}$ & $0.937$ & $0.966$ & $+0.029$ & $0.051$ & $0.914$ & $0.953$ & $+0.039$\\
$S$ & $0.863$ & $0.918$ & $+0.055$ & $0.188$ & $0.812$ & $0.881$ & $+0.069$\\
$\omega$ & $0.794$ & $0.881$ & $+0.088$ & $0.246$ & $0.730$ & $0.853$ & $+0.123$\\
\midrule
\multicolumn{8}{l}{\emph{near} $LO(0,0,1)$}\\
$T_B$ & $0.828$ & $0.835$ & $+0.007$ & $0.161$ & $0.771$ & $0.764$ & $-0.007$\\
$\Delta H_f$ & $0.832$ & $0.842$ & $+0.010$ & $0.152$ & $0.800$ & $0.807$ & $+0.006$\\
$\Delta H_{\mathrm{vap}}$ & $0.966$ & $0.969$ & $+0.003$ & $0.042$ & $0.956$ & $0.957$ & $+0.002$\\
$S$ & $0.931$ & $0.948$ & $+0.017$ & $0.045$ & $0.892$ & $0.920$ & $+0.028$\\
$\omega$ & $0.929$ & $0.945$ & $+0.016$ & $0.074$ & $0.907$ & $0.926$ & $+0.019$\\
\bottomrule
\end{tabular}
\end{table}

Three patterns emerge. First, every in-sample gain
$\Delta_{\mathrm{in}}$ is smaller than the corresponding
single-descriptor bootstrap uncertainty scale; this comparison is
descriptive and is not a formal paired or post-selection test of the
gain. Second, under the foldwise nested evaluation (within the globally
selected anchor pools) these descriptive point-estimate gains do not vanish --- tuning near $M_2$ raises the out-of-sample
correlation for $T_B$ and $\Delta H_f$ by $+0.13$ and $+0.15$, and
tuning near ${}^{m}\!M_2$ raises $\omega$ by $+0.12$ --- but in every
such case the anchor was a \emph{weak} descriptor for that property, and
the tuned correlation only climbs part way toward the level a strong
reduction already attains (tuning near $M_2$ lifts $T_B$ to $0.46$,
whereas ${}^{m}\!M_2$ already reaches $0.78$). Where the anchor is
already the column winner --- $M_2$ on $\omega$, $HM$ on $S$ --- the
nested gain is zero or slightly negative. Third, tuning near the
pure-$\gamma$ points $LO(0, 0, 1)$ and $LO(0, 0, 2)$ does not exceed the best fixed
point estimate in the tested comparison set on four of the five
properties (Table~\ref{tab:lo_tuning_bestfixed}). On
$\Delta H_{\mathrm{vap}}$ the tuned value exceeds the fixed benchmark
$LO(0, 0, 2)$ by $0.002$ in sample ($0.969$ vs $0.967$) and by $0.003$
under leave-one-out ($0.959$ vs $0.956$); these small post-selection
differences do not establish superiority. On $T_B$, $\Delta H_f$, $S$,
and $\omega$ the best fixed reduction is higher both in and out of
sample.

\begin{table}[H]
\centering
\caption{Decimal tuning near the pure-$\gamma$ points $LO(0,0,1)$ and $LO(0,0,2)$ versus the best fixed index in the tested comparison set for each octane property. The best fixed index is ranked by in-sample $|r|$ (its leave-one-out value in parentheses); the decimal columns give the best of $50$ random triples within $\pm 0.3$ of either pure-$\gamma$ point, in sample and under foldwise nested leave-one-out within the globally selected pools. Indices are ranked by in-sample $|r|$ rather than $\mathrm{LOO}$, as the latter is inflated by degenerate fits for near-zero-correlation indices. Tuning does not exceed the best fixed point estimate on four properties; on $\Delta H_{\mathrm{vap}}$ it exceeds it by $0.002$ in sample and $0.003$ under LOO, small post-selection differences that do not establish superiority.}
\label{tab:lo_tuning_bestfixed}
\small
\begin{tabular}{lcrrrr}
\toprule
Property & best fixed & $|r|_{\mathrm{fix}}$ & $\mathrm{LOO}_{\mathrm{fix}}$ & $|r|_{\mathrm{dec}}$ & $\mathrm{LOO}_{\mathrm{dec}}$\\
\midrule
$T_B$ & ${}^{m}\!M_2$ & $0.855$ & $0.781$ & $0.834$ & $0.763$\\
$\Delta H_f$ & ${}^{m}\!M_2$ & $0.891$ & $0.865$ & $0.850$ & $0.809$\\
$\Delta H_{\mathrm{vap}}$ & $LO(0,0,2)$ & $0.967$ & $0.956$ & $0.969$ & $0.959$\\
$S$ & $HM$ & $0.975$ & $0.967$ & $0.949$ & $0.921$\\
$\omega$ & $M_2$ & $0.988$ & $0.985$ & $0.946$ & $0.929$\\
\bottomrule
\end{tabular}
\end{table}

In short, local tuning raised several nested leave-one-out point
estimates for weak anchors, and except for the small
$\Delta H_{\mathrm{vap}}$ margins noted above did not exceed the best
fixed benchmark in the tested comparison set; no formal superiority
claim is made. All
tabled nested values correspond to positive signed
prediction--observation correlations with positive predictive $Q^{2}$,
and a constant-descriptor control confirms that absolute pooled
correlation alone can be inflated for degenerate descriptors (cf.\ the
caveat in the caption of Table~\ref{tab:lo_tuning_bestfixed}).

\subsection{Does the imbalance parameter add predictive information?
A matched-budget nested comparison with the parent
family}\label{sec:ablation}

The analyses above tune around fixed anchors and are descriptive.
The sharper question raised by the Gutman--Milovanovi\'c framing is
direct: given identical search budgets, does freeing $\gamma$ improve
out-of-sample prediction beyond an optimized parent
$M_{\alpha,\beta}$? We answer it with a fully nested leave-one-out
design. For each property and each outer fold (one isomer held out),
$200$ candidate parameter vectors are drawn once from
$\mathrm{Uniform}(-2,2)$ per coordinate (triples
$(\alpha,\beta,\gamma)$ for the Loyola model, pairs $(\alpha,\beta)$
with $\gamma = 0$ for the parent; identical budget, fixed random
seed); on the $17$ training isomers only, every candidate is scored
by \emph{inner} leave-one-out root-mean-square error of the
one-descriptor linear regression, the best candidate is selected, the
regression is refit on all $17$, and the held-out isomer is
predicted. No quantity computed on the held-out isomer enters
selection or fitting at any point.

\begin{table}[H]
\centering
\caption{Matched-budget nested leave-one-out comparison of the
optimized parent $M_{\alpha,\beta}$ (GM) and the optimized Loyola
family ($\gamma$ free) on the five octane properties. $Q^2$ is the
outer predictive coefficient, RMSE the outer root-mean-square error;
the last column counts outer folds (of $18$) in which the Loyola
model had the strictly smaller absolute error. All selection is
repeated inside every outer fold.}
\label{tab:ablation}
\small
\begin{tabular}{lrrrrrr}
\toprule
Property & $Q^2_{\mathrm{GM}}$ & $Q^2_{LO}$ & $\Delta Q^2$ & $\mathrm{RMSE}_{\mathrm{GM}}$ & $\mathrm{RMSE}_{LO}$ & LO better\\
\midrule
$T_B$ & $+0.613$ & $+0.748$ & $+0.135$ & $3.700$ & $2.983$ & $10/18$\\
$\Delta H_f$ & $+0.855$ & $+0.549$ & $-0.306$ & $0.463$ & $0.816$ & $4/18$\\
$\Delta H_{\mathrm{vap}}$ & $+0.859$ & $+0.935$ & $+0.076$ & $0.137$ & $0.093$ & $12/18$\\
$S$ & $+0.915$ & $+0.877$ & $-0.038$ & $1.310$ & $1.573$ & $4/18$\\
$\omega$ & $+0.979$ & $+0.957$ & $-0.022$ & $0.005$ & $0.007$ & $5/18$\\
\bottomrule
\end{tabular}
\end{table}

The answer is mixed and, on balance, negative
(Table~\ref{tab:ablation}): freeing $\gamma$ improves nested
predictive performance on two of the five properties ($T_B$,
$\Delta Q^2 = +0.135$, and $\Delta H_{\mathrm{vap}}$,
$\Delta Q^2 = +0.076$) and degrades it on three ($\Delta H_f$, $S$,
$\omega$), with the $\Delta H_f$ degradation driven by
selection instability: the parent model selected the same
$(\alpha,\beta)$ in all $18$ folds there, while the three-parameter
search selected widely varying triples (fold-to-fold standard
deviation above $1.2$ in both $\alpha$ and $\beta$). The selected
$\gamma$ values are always interior to the search box (maximum
$|\gamma| = 0.47$, never within $0.5$ of the boundary), so the
mixed outcome is not a range artefact. On the one property where
$\gamma$ helps most stably, $\Delta H_{\mathrm{vap}}$, the identical
triple $(-0.27, +0.27, +0.24)$ was selected in every fold. With
$n = 18$ molecules these $Q^2$ differences carry substantial
selection noise and we attach no inferential claim to them.

Because the comparison depends on a random candidate stream, the full
nested design was repeated for one hundred seeds ($0$--$99$) at both
budgets ($200$ and $500$ candidates per model; $1000$ configurations
in all), together with a paired control in which the Loyola
candidates are compared against the \emph{same}
$(\alpha,\beta)$ coordinates with $\gamma$ set to zero, isolating the
effect of $\gamma$ from the independent candidate streams. Freeing
$\gamma$ \emph{usually} improved nested predictive performance for
$\Delta H_{\mathrm{vap}}$ ($91/100$ seeds at budget $200$, $93/100$
at budget $500$; the paired control gives $93/100$ at both budgets)
and \emph{usually} degraded it for $\Delta H_f$ ($86/100$ and
$76/100$ seeds), but neither direction is sign-invariant across
random candidate streams; $T_B$ is close to an even split at
budget $200$ ($53/100$) but leans towards improvement at budget $500$
($64/100$), and $S$ and $\omega$ usually degrade slightly. These are descriptive
tendencies under the tested random-search procedure, not inferential
evidence of predictive superiority or degradation. The honest
summary is therefore: at matched budget, the imbalance parameter
tends to help on one property ($\Delta H_{\mathrm{vap}}$), tends to
hurt on three ($\Delta H_f$, $S$, $\omega$), and is budget-dependent on
$T_B$, with no sign-invariant effect in either
direction. Its value lies in the structural properties established
in Section~\ref{sec:geometry} and the discrimination behaviour
below, not in systematic predictive gains.

\subsection{Degeneracy across tree orders}\label{sec:degeneracy}

Degeneracy measures an index's ability to discriminate
non-isomorphic graphs~\cite{Konstantinova1996Degeneracy}; lower is
better. On the $106$ non-isomorphic
trees of order $10$ we use the distinct-value convention
\[
\mathrm{Degeneracy}(\%) = 100\Bigl(1 - \tfrac{k}{N}\Bigr),
\quad k = \#\{\text{distinct index values}\},\ N = 106,
\]
i.e.\ the fraction of the $N$ structures that fail to receive a
unique value. The pure-$\gamma$ points $LO(0, 0, 1)$ and
$LO(0, 0, 2)$ attain degeneracy $25.5\%$ ($79$ distinct values),
among the lowest of all compared indices and tied with $R$, $GA$,
and $AG$; indeed the $106$ trees realise only $79$ distinct
edge-degree-pair count vectors, so $25.5\%$ is the minimum degeneracy
attainable by \emph{any} degree-pair-based index on this set, and the
pure-$\gamma$ points attain this floor --- for vertex-degree-based
weights $\varphi_{ij}$ that are algebraic numbers,
Rada~\cite{Rada2019ExpVDB} proved that the exponential index
$e^{\varphi}$ discriminates the equivalence classes of the
corresponding $\widehat{m}_{\varphi}$-structure relation; the
pure-$\gamma$ imbalance weights used here are rational, hence
algebraic, so that theorem separates the classes of graphs with
different $q$-value histograms (not the degree-pair profiles
themselves, since $q$ is not injective on pairs: $q_{12} = q_{24}$
and $q_{ii} = 0$), and attainment of the
$79$-profile floor on this particular order-$10$ tree set is the
separate computational observation reported here; the first Zagreb index $LO(0, 1, 0)$ ($83.0\%$) and the
second Zagreb index $LO(1, 0, 0)$ ($67.0\%$) are the most degenerate
(Figure~\ref{fig:degeneracy}).

\begin{figure}[H]
\centering
\safefig{fig_lo_degeneracy_order10_trees_v2.pdf}{0.9}
\caption{Degeneracy percentage of the compared indices on the $106$ non-isomorphic trees of order $10$.}
\label{fig:degeneracy}
\end{figure}

To test whether this floor attainment is peculiar to order $10$, the
same computation was repeated for all non-isomorphic trees of orders
$7$ through $12$ and, separately, for the molecular subsets with
$\Delta \le 4$ (Table~\ref{tab:multiorder}). The pure-$\gamma$ points
attain the degree-pair-profile floor at \emph{every} order in both
sets. Honesty requires noting that the Randi\'c index also attains
the floor at every order in this range (while $\chi$ falls off the
floor from order $10$), so floor attainment distinguishes the
pure-$\gamma$ points from the most degenerate classical indices
($M_1$, $M_2$) but not from $R$; the observation is that the
imbalance weights lose nothing, not that they are unique in this
respect. Attainment on orders $7$--$12$ is a computational
observation, not a proof of injectivity on degree-pair profiles for
all orders. Indeed attainment fails at larger orders: because
$LO(\cdot\,; 0, 0, \gamma)$ depends only on the histogram of
$q$-values, two trees with distinct degree-pair profiles but equal
$q$-histograms receive equal values for every $\gamma$. Among the
trees of order $13$ the profiles
$\{(1,2)^{1}, (1,3)^{2}, (1,6)^{5}, (2,3)^{2}, (2,6)^{1}, (3,3)^{1}\}$
and $\{(1,3)^{3}, (1,6)^{5}, (2,2)^{1}, (2,3)^{2}, (3,6)^{1}\}$
(exponents are edge multiplicities) share the $q$-histogram
$\{0^{1}, \tfrac15^{2}, \tfrac13^{1}, \tfrac12^{3}, \tfrac57^{5}\}$,
and the pure-$\gamma$ points realise $566$ distinct values against
$570$ profiles; among molecular trees ($\Delta \le 4$) the first
failure occurs at order $16$ ($809$ values against $810$ profiles).

\begin{table}[H]
\centering
\caption{Discrimination across tree orders $n = 7$--$12$: number of
trees $N$, number of distinct edge-degree-pair profiles (the floor
attainable by any vertex-degree-based index), and distinct-value
counts for representative indices. ``mol.''\ denotes the molecular
subset $\Delta \le 4$.}
\label{tab:multiorder}
\small
\begin{tabular}{llrrrrrrr}
\toprule
$n$ & set & $N$ & floor & $M_1$ & $M_2$ & $R$ & $LO(0,0,1)$ & $LO(0,0,2)$\\
\midrule
$7$ & all & $11$ & $11$ & $7$ & $9$ & $11$ & $11$ & $11$\\
$7$ & mol. & $9$ & $9$ & $5$ & $7$ & $9$ & $9$ & $9$\\
$8$ & all & $23$ & $21$ & $9$ & $17$ & $21$ & $21$ & $21$\\
$8$ & mol. & $18$ & $16$ & $6$ & $13$ & $16$ & $16$ & $16$\\
$9$ & all & $47$ & $40$ & $13$ & $22$ & $40$ & $40$ & $40$\\
$9$ & mol. & $35$ & $28$ & $7$ & $16$ & $28$ & $28$ & $28$\\
$10$ & all & $106$ & $79$ & $18$ & $35$ & $79$ & $79$ & $79$\\
$10$ & mol. & $75$ & $49$ & $8$ & $21$ & $49$ & $49$ & $49$\\
$11$ & all & $235$ & $152$ & $21$ & $43$ & $152$ & $152$ & $152$\\
$11$ & mol. & $159$ & $83$ & $9$ & $25$ & $83$ & $83$ & $83$\\
$12$ & all & $551$ & $294$ & $27$ & $57$ & $294$ & $294$ & $294$\\
$12$ & mol. & $355$ & $137$ & $10$ & $28$ & $137$ & $137$ & $137$\\
\bottomrule
\end{tabular}
\end{table}

\subsection{Structure sensitivity on decane isomers}\label{sec:ss}

The structure-sensitivity criterion for descriptor quality was
introduced by Furtula, Gutman and
Dehmer~\cite{FurtulaGutmanDehmer2013SS}, with a fingerprint-based
variant due to Raki\'c and Furtula~\cite{RakicFurtula2019SS}. We
implement the published protocol on the $75$ decane isomer trees
$\Psi$: for each tree $G$, the similar-structure set
$S(G) \subset \Psi$ consists of the non-isomorphic trees obtainable
from $G$ by relocating a single edge within the molecular class
(graph edit distance two, $\Delta \le 4$ preserved); then
\begin{gather*}
SS(TI, G) = \frac{1}{|S(G)|} \sum_{\Omega \in S(G)}
\frac{|TI(\Omega) - TI(G)|}{TI(G)},\\
Abr(TI, G) = \max_{\Omega \in S(G)}
\frac{|TI(\Omega) - TI(G)|}{TI(G)},
\end{gather*}
and the class-level $SS(TI)$ and $Abr(TI)$ are the means over all
$75$ trees. On this class every $S(G)$ is nonempty (sizes $7$ to
$43$), the neighbour relation is symmetric, and $Abr \ge SS$ holds
throughout, as it must. As an external control, the implementation
reproduces the published decane-class values of Barman and
Das~\cite{BarmanDas2026HSO} for the seven indices common to both
studies ($M_1$, ${}^{m}\!M_2$, $R$, $\chi$, $H$, $GA$, $AG$;
fourteen tabulated $SS$/$Abr$ values): eleven agree exactly at the four
displayed decimals and the remaining three differ by one unit in the
fourth decimal. Under the published desideratum ($SS$ large,
$Abr$ small~\cite{FurtulaGutmanDehmer2013SS}), the pure-$\gamma$
point $LO(0,0,2)$ attains the largest structure sensitivity in the
comparison set ($SS = 0.1257$, marginally above the hyper-Zagreb
index at $0.1256$) with a markedly smaller abruptness than
hyper-Zagreb ($0.274$ versus $0.313$), and the two pure-$\gamma$
points have the two largest sensitivity-to-abruptness ratios
($0.4579$ and $0.4565$). Given the redundancy analysis of
Section~\ref{sec:octanes}, the $SS$ margin over hyper-Zagreb should
be read as a tie in sensitivity at smaller abruptness, not as a
ranking.

\begin{table}[H]
\centering
\caption{Structure sensitivity $SS$, abruptness $Abr$, and their
ratio on the $75$ decane isomers under the published
similar-structure protocol of~\cite{FurtulaGutmanDehmer2013SS}
(single-edge relocation within the molecular class).}
\label{tab:fgdss}
\small
\begin{tabular}{lrrr}
\toprule
Index & $SS$ & $Abr$ & $SS/Abr$\\
\midrule
$M_1$ & $0.0566$ & $0.1336$ & $0.4238$\\
$M_2$ & $0.0922$ & $0.2443$ & $0.3776$\\
$HM$ & $0.1256$ & $0.3126$ & $0.4017$\\
${}^{m}\!M_2$ & $0.0516$ & $0.1177$ & $0.4384$\\
$R$ & $0.0264$ & $0.0588$ & $0.4481$\\
$\chi$ & $0.0264$ & $0.0591$ & $0.4463$\\
$H/2$ & $0.0516$ & $0.1143$ & $0.4514$\\
$ISI$ & $0.0226$ & $0.0597$ & $0.3795$\\
$GA$ & $0.0241$ & $0.0535$ & $0.4507$\\
$AG$ & $0.0256$ & $0.0579$ & $0.4415$\\
$LO(0,0,1)$ & $0.0624$ & $0.1367$ & $0.4565$\\
$LO(0,0,2)$ & $0.1257$ & $0.2745$ & $0.4579$\\
\bottomrule
\end{tabular}
\end{table}

Earlier drafts of this study reported the coefficient of variation
$\sigma/\mu$ and relative range $(\max-\min)/\mu$ over the isomer
class under the names $SS$ and $Abr$; those are global dispersion
summaries, not the published pairwise quantities, and they are
retained only as internal artefacts of the reproducibility package
under the names \emph{normalized dispersion} and \emph{normalized
range}.

% =====================================================================
\section{Conclusion and further directions}\label{sec:conclusion}
% =====================================================================

The Loyola family is the normalized degree-imbalance exponential
extension of the Gutman--Milovanovi\'c family
$M_{\alpha,\beta}$~\cite{GutmanMilovanovic2020AKCE,
MolinaSigarreta2025GM}, to which it reduces identically at
$\gamma = 0$. What the third parameter contributes is now stated
precisely: the imbalance direction is affinely independent of the two
parent directions on the admissible degree-pair set for $\Delta \ge 3$
(Proposition~\ref{prop:indep}; a class-level statement, see
Remark~\ref{rem:ident}); the family has a globally convex
logarithm whose $\gamma$-derivatives are the mean and variance of the
edge imbalance under an induced edge measure
(Theorem~\ref{thm:logconvex}); and the two-sided bounds of
Section~\ref{sec:bounds} extend inequalities known for the parent to
the $\gamma \neq 0$ regime, several of them sharp.

Chemometrically the conclusions are deliberately honest. Every
vertex-degree-based index on the octane benchmark operates above an
irreducible floor set by two degree-pair collision pairs
(Proposition~\ref{prop:collision}); the twelve-index comparison set
is effectively two-dimensional, so third-decimal correlation
differences carry no ranking information; and at matched search
budget, across one hundred random-search seeds and two budgets, the
free imbalance parameter usually improves nested out-of-sample
prediction on one of five properties, usually degrades it on three,
and is mixed on the fifth, with neither tendency sign-invariant across candidate streams
(Table~\ref{tab:ablation}) --- it is not a systematic predictive
upgrade over the optimized parent. Its measurable contribution is
structural: the pure-$\gamma$ points attain the degree-pair
discrimination floor at every tree order $7$--$12$ tested, and under
the published structure-sensitivity protocol they combine
hyper-Zagreb-level sensitivity with smaller abruptness and the
largest sensitivity-to-abruptness ratios in the comparison set. The
Loyola family is hence best viewed as a provably independent,
analytically transparent third direction on the
Gutman--Milovanovi\'c parameter space, confined --- like every
vertex-degree-based family --- to the finite-dimensional degree-pair
descriptor space on molecular graphs.

Building on the Loyola index, the following open problems indicate natural directions for further work:
\begin{enumerate}
\item Sharp two-sided bounds for $\Wuzi$ via degree-based indices
lying outside the family, such as the atom-bond-connectivity index $ABC$ and the
Sombor index.


\item Extremal-graph characterisations for several graph classes, including trees, unicyclic and
bicyclic graphs.

\item Chemometric extension to additional chemical compound classes such as benzene derivatives, alkane homologues beyond octane and
decane and to compound-class-specific QSPR endpoints.
\end{enumerate}

\acknowledgment{The authors thank colleagues for helpful discussions
of the computational extremal-graph search.}

\appendix


\section{Octane isomer property data}\label{app:data}

For reproducibility, Table~\ref{tab:octane-data} lists the full
$18$-isomer dataset used in Section~\ref{sec:sensitivity}. Units:
$T_B$ in $^{\circ}$C; $\Delta H_f$ and $\Delta H_{\mathrm{vap}}$ in
kcal\,mol$^{-1}$ ($\Delta H_{\mathrm{vap}}$ at $298$~K); $S$ in
cal\,mol$^{-1}$K$^{-1}$; $\omega$ dimensionless.

\begin{table}[H]
\centering\small
\caption{Physicochemical properties of the $18$ octane isomers.}
\label{tab:octane-data}
\begin{tabular}{lrrrrr}
\toprule
Isomer & $T_B$ & $\Delta H_f$ & $\Delta H_{\mathrm{vap}}$ & $S$ & $\omega$\\
\midrule
$n$-octane & 125.6 & -49.82 & 9.92 & 111.55 & 0.398\\
2-methylheptane & 117.6 & -51.5 & 9.48 & 109.84 & 0.378\\
3-methylheptane & 118.9 & -50.82 & 9.52 & 111.26 & 0.371\\
4-methylheptane & 117.7 & -50.69 & 9.48 & 109.32 & 0.372\\
2,2-dimethylhexane & 106.8 & -53.71 & 8.92 & 103.13 & 0.339\\
2,3-dimethylhexane & 115.6 & -51.13 & 9.27 & 108.02 & 0.348\\
2,4-dimethylhexane & 109.4 & -52.44 & 9.03 & 106.98 & 0.344\\
2,5-dimethylhexane & 109.1 & -53.21 & 9.05 & 105.72 & 0.357\\
3,3-dimethylhexane & 111.9 & -52.61 & 9.04 & 104.74 & 0.322\\
3,4-dimethylhexane & 117.7 & -50.91 & 9.32 & 106.59 & 0.34\\
2-methyl-3-ethylpentane & 115.6 & -50.48 & 9.21 & 106.06 & 0.33\\
3-ethyl-3-methylpentane & 118.3 & -51.38 & 9.21 & 101.48 & 0.302\\
3-ethylhexane & 118.5 & -50.4 & 9.48 & 109.43 & 0.362\\
2,2,3-trimethylpentane & 109.8 & -52.61 & 8.88 & 101.31 & 0.3\\
2,2,4-trimethylpentane & 99.2 & -53.57 & 8.4 & 101.81 & 0.305\\
2,3,3-trimethylpentane & 114.8 & -51.73 & 9.02 & 101.31 & 0.291\\
2,3,4-trimethylpentane & 113.5 & -51.97 & 9.01 & 102.39 & 0.317\\
2,2,3,3-tetramethylbutane & 106.5 & -53.99 & 8.41 & 93.06 & 0.247\\
\bottomrule
\end{tabular}
\end{table}



% References should be ordered alphabetically.
\singlespacing

 \begin{thebibliography}{99}

  \bibitem{AhmadFarooqDas2025Extremal}
  S. Ahmad, R. Farooq, K. C. Das, The general Sombor index of extremal
  trees with a given maximum degree, \textit{MATCH Commun. Math. Comput.
  Chem.\/} \textbf{94} (2025) 825--853.

  \bibitem{Albertson1997}
  M. O. Albertson, The irregularity of a graph, \textit{Ars Combin.\/}
  \textbf{46} (1997) 219--225.

  \bibitem{ArockiarajPaulClement2023Hybrid} M. Arockiaraj, D. Paul, J.
  Clement, S. Tigga, K. Jacob, K. Balasubramanian, Novel molecular hybrid
  geometric-harmonic-Zagreb degree based descriptors and their efficacy in
  QSPR studies of polycyclic aromatic hydrocarbons, \textit{SAR QSAR
  Environ. Res.\/} \textbf{34} (2023) 569--589.

  \bibitem{BarmanDas2026HSO} J. Barman, S. Das, Geometric approach to
  degree-based topological index: Hyperbolic Sombor index, \textit{MATCH
  Commun. Math. Comput. Chem.\/} \textbf{95} (2026) 63--94.

  \bibitem{BollobasErdos1998}
  B. Bollob\'as, P. Erd\H{o}s, Graphs of extremal weights, \textit{Ars Combin.\/}
  \textbf{50} (1998) 225--233.

  \bibitem{DasCevikCangulShang2021Sombor} K. C. Das, A. S. \c{C}evik, I. N.
  Cangul, Y. Shang, On Sombor index, \textit{Symmetry\/} \textbf{13}
  (2021) \#140.

  \bibitem{DasGutman2022SomborTrees} K. C. Das, I. Gutman, On Sombor index
  of trees, \textit{Appl. Math. Comput.\/} \textbf{412} (2022) \#126575.

  \bibitem{FurtulaGutman2015F} B. Furtula, I. Gutman, A forgotten
  topological index, \textit{J. Math. Chem.\/} \textbf{53} (2015)
  1184--1190.

  \bibitem{FurtulaGutmanDehmer2013SS} B. Furtula, I. Gutman, M. Dehmer, On
  structure-sensitivity of degree-based topological indices, \textit{Appl.
  Math. Comput.\/} \textbf{219} (2013) 8973--8978.

  \bibitem{GaoAMC2024} W. Gao, Y. Gao, The extremal trees for
  exponential vertex-degree-based topological indices, \textit{Appl.
  Math. Comput.\/} \textbf{472} (2024) 128634.

  \bibitem{GranadosJMC2025GM} A. Granados, A. Portilla, Y. Quintana,
  E. Tour\'is, Bounds for the Gutman--Milovanovi\'c index and some
  applications, \textit{J. Math. Chem.\/} \textbf{63} (2025)
  406--418.

  \bibitem{GujiWali2025Lanzhou} R. Guji, M. Wali, The Sombor index
  (coindex) and Lanzhou index (coindex) of some graphs, \textit{Axioms\/}
  \textbf{14} (2025) 164.

  \bibitem{GutmanMilovanovic2020AKCE} I. Gutman, E. Milovanovi\'c,
  I. Milovanovi\'c, Beyond the Zagreb indices, \textit{AKCE Int. J.
  Graphs Comb.\/} \textbf{17} (2020) 74--85.

  \bibitem{GutmanToganYurttasCevikCangul2018Sigma} I. Gutman, M. Togan, A.
  Yurtta\c{s}, A. S. \c{C}evik, I. N. Cangul, Inverse problem for sigma index,
  \textit{MATCH Commun. Math. Comput. Chem.\/} \textbf{79} (2018) 491--508.

  \bibitem{Gutman2021Sombor} I. Gutman, Geometric approach to degree-based
  topological indices: Sombor indices, \textit{MATCH Commun. Math. Comput.
  Chem.\/} \textbf{86} (2021) 11--16.

  \bibitem{Gutman2021SomborBasic} I. Gutman, Some basic properties of
  Sombor indices, \textit{Open J. Discr. Appl. Math.\/} \textbf{4} (2021)
  1--3.

  \bibitem{GutmanFurtulaOz2024Elliptic} I. Gutman, B. Furtula, M. S. Oz,
  Geometric approach to vertex-degree-based topological indices: Elliptic
  Sombor index, theory and application, \textit{Int. J. Quantum Chem.\/}
  \textbf{124} (2024) \#e27346.

  \bibitem{Gutman2026SomborSurvey} I. Gutman, Survey of Sombor-type
  degree-based topological indices, \textit{MATCH Commun. Math. Comput.
  Chem.\/} \textbf{96} (2026) 819--837.

  \bibitem{HuMATCH2022} Z. Hu, L. Li, X. Li, D. Peng, Extremal graphs
  for topological index defined by a degree-based edge-weight
  function, \textit{MATCH Commun. Math. Comput. Chem.\/} \textbf{88}
  (2022) 505--520.

  \bibitem{KalaamGreeniArockiaraj2024} A. R. A. Kalaam, A. B. Greeni, M.
  Arockiaraj, Modified reverse degree descriptors for combined topological
  and entropy characterizations of 2D metal organic frameworks:
  applications in graph energy prediction, \textit{Front. Chem.\/}
  \textbf{12} (2024) \#1470231.

  \bibitem{Konstantinova1996Degeneracy} E. V. Konstantinova, The
  discrimination ability of some topological and information distance indices
  for graphs of unbranched hexagonal systems, \textit{J. Chem. Inf. Comput.
  Sci.\/} \textbf{36} (1996) 54--57.

  \bibitem{KosariDehgardiKhan2023Lower} S. Kosari, N. Dehgardi, A. Khan,
  Lower bound on the KG-Sombor index, \textit{Commun. Comb. Optim.\/}
  \textbf{8} (2023) 751--757.


  \bibitem{NISTWebBook} P. J. Linstrom, W. G. Mallard (Eds.), \textit{NIST
  Chemistry WebBook, NIST Standard Reference Database Number 69}, National
  Institute of Standards and Technology, Gaithersburg MD,
  \href{https://doi.org/10.18434/T4D303}{doi:10.18434/T4D303} (retrieved January 2025).

  \bibitem{MendezBermudezAguilarSanchez2022MeanSombor} J. A.
  M\'endez-Berm\'udez, R. Aguilar-S\'anchez, E. D. Molina, J. M. Rodr\'iguez, Mean
  Sombor index, \textit{Discr. Math. Lett.\/} \textbf{9} (2022) 18--25.

  \bibitem{MilovanovicMilovanovicMatejic2021Sombor} I. Milovanovi\'c, E.
  Milovanovi\'c, M. Mateji\'c, On some mathematical properties of Sombor
  indices, \textit{Bull. Int. Math. Virtual Inst.\/} \textbf{11} (2021)
  341--353.

  \bibitem{MolinaSigarreta2025GM} E. D. Molina, J. M.
  Rodr\'iguez-Garc\'ia, J. M. Sigarreta, S. J. Torralbas Fitz, On the
  Gutman--Milovanovi\'c index and chemical applications, \textit{AIMS
  Math.\/} \textbf{10} (2025) 1998--2020.

  \bibitem{Movahedi2025DSO} F. Movahedi, Diminished Sombor index and its
  relationship with topological indices, \textit{Filomat\/} \textbf{39}
  (2025) 10519--10532.

  \bibitem{MovahediGutmanRedzepovicFurtula2026} F. Movahedi, I. Gutman, I.
  Red\v{z}epovi\'c, B. Furtula, Diminished Sombor index, \textit{MATCH Commun.
  Math. Comput. Chem.\/} \textbf{95} (2026) 141--162.

  \bibitem{PortillaJMC2025GM} A. Portilla, J. M. Rodr\'iguez,
  I. Salas Lorenzo, J. M. Sigarreta, The Gutman--Milovanovi\'c index
  in mathematical chemistry, \textit{J. Math. Chem.\/} \textbf{63}
  (2025) 2068--2087.

  \bibitem{Rada2019ExpVDB} J. Rada, Exponential vertex-degree-based
  topological indices and discrimination, \textit{MATCH Commun. Math.
  Comput. Chem.\/} \textbf{82} (2019) 29--41.

  \bibitem{Rada2026LinearVDB} J. Rada, A finite-dimensional linear
  framework for vertex-degree-based topological indices, \textit{MATCH
  Commun. Math. Comput. Chem.\/} \textbf{96} (2026) 841--863.

  \bibitem{RakicFurtula2019SS} M. Raki\'c, B. Furtula, A novel method for
  measuring the structure sensitivity of molecular descriptors, \textit{J.
  Chemometrics\/} \textbf{33} (2019) e3138.

  \bibitem{Shang2022Simplicial}
  Y. Shang, Sombor index and degree-related properties of simplicial
  networks, \textit{Appl. Math. Comput.\/} \textbf{419} (2022) 126881.

  \bibitem{Sigarreta2022ExpVDB} J. M. Sigarreta, Extremal problems on
  exponential vertex-degree-based topological indices, \textit{Math.
  Biosci. Eng.\/} \textbf{19} (2022) 6985--6995.

  \bibitem{VukicevicFurtula2009GA} D. Vuki\v{c}evi\'c, B. Furtula, Topological
  index based on the ratios of geometrical and arithmetical means of
  end-vertex degrees of edges, \textit{J. Math. Chem.\/} \textbf{46} (2009)
  1369--1376.

  \bibitem{WangMaoLiFurtula2022SomborRelations} Z. Wang, Y. Mao, Y. Li, B.
  Furtula, On relations between Sombor and other degree-based indices,
  \textit{J. Appl. Math. Comput.\/} \textbf{68} (2022) 1--17.

  \bibitem{ZhouTrinajstic2009sumconn}
  B. Zhou, N. Trinajsti\'c, On a novel connectivity index, \textit{J. Math.
  Chem.\/} \textbf{46} (2009) 1252--1270.

  \bibitem{ZhouTrinajstic2010} B. Zhou, N. Trinajsti\'c, On general
  sum-connectivity index, \textit{J. Math. Chem.\/} \textbf{47} (2010)
  210--218.

  \bibitem{ZuoRatherImranAli2022Modified}
  X. Zuo, B. A. Rather, M. Imran, A. Ali, On some topological indices
  defined via the modified Sombor matrix, \textit{Molecules\/} \textbf{27}
  (2022) 6772.

  \end{thebibliography}

\end{document}
